\documentclass[11pt,a4paper]{article}
\usepackage[a4paper,margin=2cm,top=1.8cm,bottom=2cm]{geometry}
\usepackage{amsmath,amssymb}
\usepackage{fontspec}
\usepackage{pifont}
\usepackage{xcolor}
\usepackage[most]{tcolorbox}
\usepackage{enumitem}
\usepackage{booktabs}
\usepackage{fancyhdr}
\usepackage{extramarks}
\usepackage{needspace}
\usepackage[hidelinks]{hyperref}

\setlength{\parindent}{0pt}
\setlength{\parskip}{3pt}

% ---------- colours (same family as the formula sheet) ----------
\definecolor{accent}{HTML}{1D4E89}
\definecolor{accentlight}{HTML}{EAF1FA}
\definecolor{muted}{HTML}{4A5560}
\definecolor{rem}{HTML}{8A5A00}
\definecolor{remlight}{HTML}{FFF7E6}
\definecolor{edge}{HTML}{B03A2E}
\definecolor{edgelight}{HTML}{FBEDEB}
\definecolor{good}{HTML}{1E7B3A}
\definecolor{goodlight}{HTML}{EAF6EE}

% ---------- symbols ----------
\newcommand{\cmark}{\textcolor{good}{\ding{51}}}
\newcommand{\xmark}{\textcolor{edge}{\ding{55}}}
\newcommand{\sthree}{\textcolor{edge}{\ding{72}\ding{72}\ding{72}}}
\newcommand{\stwo}{\textcolor{rem}{\ding{72}\ding{72}}}
\newcommand{\sone}{\textcolor{muted}{\ding{72}}}

% ---------- header / footer ----------
\pagestyle{fancy}\fancyhf{}
\fancyhead[L]{\small\color{muted}MXB103 Derivations}
\fancyhead[R]{\small\color{muted}\ifx\topleftxmark\empty\firstleftxmark\else\topleftxmark\fi}
\fancyfoot[C]{\small\color{muted}\thepage}
\renewcommand{\headrulewidth}{0.3pt}

% ---------- section headings ----------
\newcommand{\topic}[1]{\clearpage\extramarks{#1}{#1}\phantomsection\addcontentsline{toc}{section}{#1}%
  {\color{accent}\Large\bfseries #1\par}\vspace{-2pt}{\color{accent}\rule{\linewidth}{1pt}}\par\vspace{6pt}}

% ---------- derivation box ----------
\newcounter{stepc}
\newtcolorbox{deriv}[3]{enhanced,breakable,colback=white,colframe=accent!65,colbacktitle=accentlight,
  coltitle=accent,fonttitle=\bfseries\large,boxrule=0.6pt,arc=2pt,left=8pt,right=8pt,top=4pt,bottom=6pt,
  title={#1\quad #2},before skip=10pt,after skip=10pt,
  before upper={\setcounter{stepc}{0}\phantomsection\addcontentsline{toc}{subsection}{#1\quad #2}%
    {\small\textbf{Priority:} #3}\par\vspace{3pt}}}

\newcommand{\goal}[1]{{\color{muted}\textbf{Goal:}}\ #1\par\vspace{2pt}}
\newcommand{\idea}[1]{\begin{tcolorbox}[enhanced,colback=remlight,colframe=rem!40,boxrule=0.4pt,arc=2pt,left=6pt,right=6pt,top=2pt,bottom=2pt]
  {\color{rem}\textbf{Idea.}}\ #1\end{tcolorbox}}
\newcommand{\step}[1]{\par\vspace{5pt}\refstepcounter{stepc}\noindent{\color{accent}\textbf{Step \thestepc.}}\ #1\par}
\newcommand{\mk}[1]{\textcolor{edge}{\textbf{[#1]}}}
\newcommand{\why}[1]{{\small\color{muted}\textit{Why:} #1\par}}
\newcommand{\trap}[1]{{\small\textcolor{edge}{\textbf{Common mistake:}}\ #1\par}}
\newcommand{\result}[1]{\begin{tcolorbox}[enhanced,colback=goodlight,colframe=good!50,boxrule=0.5pt,arc=2pt,left=6pt,right=6pt,top=3pt,bottom=3pt]
  {\color{good}\textbf{Result.}}\ #1\end{tcolorbox}}
\newcommand{\chk}[1]{{\small\textcolor{good}{\textbf{Check:}}\ #1\par}}

\allowdisplaybreaks

\begin{document}

% ======================= TITLE & HOW TO USE =======================
\thispagestyle{plain}
{\color{accent}\huge\bfseries MXB103 Derivations\par}
\vspace{2pt}
{\Large\color{muted} Every formula in the unit, derived slowly, one step at a time\par}
\vspace{10pt}

Every derivation in this document has the same layout:

\begin{itemize}[leftmargin=16pt,itemsep=2pt]
\item \textbf{Goal}: the formula we are about to prove.
\item \textbf{Idea}: the one-sentence plan, so you know where the steps are heading.
\item \textbf{Steps}: each one says \emph{what} to do and \emph{why}. When the step matches an official marking scheme, its marks are shown in \mk{red}.
\item \textbf{Result}: the final formula in a green box.
\end{itemize}

\textbf{Priority stars}
\begin{center}
\begin{tabular}{ll}
\toprule
\sthree & Asked in the 2023, 2024 or 2025 exam. \textbf{Learn to write it from memory.}\\
\stwo & A core building block the exam relies on, or a likely variation.\\
\sone & For understanding. Not asked directly, but it explains \emph{why} a method works.\\
\bottomrule
\end{tabular}
\end{center}

\textbf{How to practise.} Cover everything below the Goal, write the derivation on blank paper, then compare step by step. A derivation is ``learnt'' when you can do it twice in a row without looking.

\vspace{6pt}
\begin{tcolorbox}[enhanced,colback=edgelight,colframe=edge!50,boxrule=0.5pt,arc=2pt,title={\sthree\ The exam derivations (learn these first)},coltitle=edge,colbacktitle=edgelight,fonttitle=\bfseries]
\small
\begin{tabular}{@{}lll@{}}
D2.1 & Cancellation-free quadratic formula & 2025 Q1b (3 marks)\\
D2.4 & Taylor polynomial + error bound (full layout) & Q2 every year (6 marks)\\
D4.1 & Total derivative $f'(t,y)$ & needed for every Taylor 2 question\\
D4.2 & Second-order Taylor method & 2023 Q5 (4 marks)\\
D4.3 & Modified Euler method & 2024 Q5, 2025 Q5 (4 marks)\\
D6.2 & Expanding $f[x_0,x_1,x_2,x_3]$ & 2023 Q8, 2025 Q8 (4 marks)\\
D7.2 & $f[x_0,x_1,x_2]=\Delta^2y_0/(2!\,h^2)$ & 2024 Q8 (4 marks)\\
D8.1 & Showing $[a,b]$ is a bracketing & Q9a every year (2 marks)\\
D9.1 & Newton's method from the tangent line & 2025 Q10 (3 marks)\\
D9.3 & Iterations for $\sqrt a$, $\sqrt[3]a$ (Newton or secant) & Q11 2024, 2025 (4 marks)\\
D10.3 & Central difference for $f''$ with error term & 2023 Q11a (4 marks)\\
D10.5 & Effect of $h$ (truncation vs roundoff) & 2023 Q11b (2 marks)\\
D11.1, D11.2 & Trapezoidal rule (one strip, then composite) & 2023 Q13 (5 marks)\\
D11.3 & Simpson's rule (one pair of strips) & 2024 Q13 (5 marks)\\
D12.3 & How LU solves $A\mathbf x=\mathbf b$ (in words) & 2023 Q15b (3 marks)\\
D12.4 & Doolittle LU formulas & Q15 every year (2--3 marks)\\
D13.1 & Positive definite by completing the square & 2025 Q16a (3 marks)\\
D13.2 & Cholesky formulas & Q16 every year (2--3 marks)\\
\end{tabular}
\end{tcolorbox}

\vspace{4pt}
{\small\tableofcontents}

% ======================= TOOLKIT =======================
\topic{Toolkit: five tools used again and again}

Most derivations in this unit are built from the same five tools. Learn these once, and the derivations become much shorter.

\begin{deriv}{T1}{Taylor's theorem written with a step $h$}{\stwo\ used in Weeks 3, 4, 9 and 10}
\goal{a way to write $f(x_0+h)$ and $f(x_0-h)$ in terms of $f$ and its derivatives at $x_0$.}

\step{Start from Taylor's theorem (Week 2) centred on $x_0$:}
\[f(x)=f(x_0)+f'(x_0)(x-x_0)+\frac{f''(x_0)}{2!}(x-x_0)^2+\cdots+\frac{f^{(n)}(x_0)}{n!}(x-x_0)^n+\frac{f^{(n+1)}(c)}{(n+1)!}(x-x_0)^{n+1}\]
\why{this is exact, with $c$ somewhere between $x_0$ and $x$.}

\step{Put $x=x_0+h$. Then $x-x_0=h$, so every bracket becomes a power of $h$:}
\[f(x_0+h)=f(x_0)+h\,f'(x_0)+\frac{h^2}{2}f''(x_0)+\frac{h^3}{6}f'''(x_0)+\cdots+\frac{h^{n+1}}{(n+1)!}f^{(n+1)}(c)\]

\step{Replace $h$ by $-h$. Even powers keep their sign and odd powers flip:}
\[f(x_0-h)=f(x_0)-h\,f'(x_0)+\frac{h^2}{2}f''(x_0)-\frac{h^3}{6}f'''(x_0)+\cdots\]

\result{$f(x_0\pm h)=f(x_0)\pm h f'(x_0)+\frac{h^2}{2}f''(x_0)\pm\frac{h^3}{6}f'''(x_0)+\cdots$\quad (signs alternate on odd powers).}
\end{deriv}

\begin{deriv}{T2}{What $O(h^p)$ means, and the rules for using it}{\stwo\ used in every error analysis}
\goal{understand why an unknown term like $\frac{y''(c)}{2}h^2$ can be written as $O(h^2)$.}

\step{A term is $O(h^p)$ if it is at most (some fixed constant) $\times\,h^p$ for small $h$.}
\why{in $\frac{y''(c)}{2}h^2$ we never know $c$, but $|y''|\le M$ on the interval, so the term is at most $\frac M2 h^2$. We keep only the part that matters: how it scales with $h$.}

\step{The rules you need (each follows from ``constant $\times\,h^p$''):}
\begin{align*}
h\cdot O(h^p)&=O(h^{p+1}) && \text{(multiplying by $h$ adds one power)}\\
\tfrac1h\cdot O(h^p)&=O(h^{p-1}) && \text{(dividing by $h$ removes one power)}\\
O(h^p)+O(h^p)&=O(h^p) && \text{(two such terms is still constant $\times h^p$)}\\
O(h^p)+O(h^q)&=O(h^p)\ \text{if }p<q && \text{(for small $h$ the lower power is bigger and dominates)}
\end{align*}

\step{What it predicts: if the error is $\approx C h^p$, halving $h$ divides the error by $2^p$.}
\result{$O(h)$: halve $h$, error $\div2$.\quad $O(h^2)$: error $\div4$.\quad $O(h^4)$: error $\div16$.}
\end{deriv}

\begin{deriv}{T3}{The IVT averaging trick (merging two unknown $c$'s into one)}{\stwo\ used in D10.2, D10.3}
\goal{if $g$ is continuous and $c_1,c_2$ are two points, show there is a single $c$ between them with $g(c)=\frac{g(c_1)+g(c_2)}{2}$.}

\step{The average of two numbers always lies between them:}
\[\min\{g(c_1),g(c_2)\}\le\frac{g(c_1)+g(c_2)}{2}\le\max\{g(c_1),g(c_2)\}\]
\step{The Intermediate Value Theorem says a continuous function takes every value between $g(c_1)$ and $g(c_2)$ somewhere between $c_1$ and $c_2$.}
\step{So it takes the average value at some point $c$ between $c_1$ and $c_2$.}
\result{$\dfrac{g(c_1)+g(c_2)}{2}=g(c)$ for some $c$ between $c_1$ and $c_2$.}
\why{finite-difference derivations produce two error terms with different unknown points. This trick tidies them into one.}
\end{deriv}

\begin{deriv}{T4}{The substitution $x=x_0+sh$ for equally spaced points}{\stwo\ used in D7.3, D11.1--D11.3}
\goal{rewrite expressions in $x$ in terms of the ``step counter'' $s$.}
\step{Equally spaced nodes satisfy $x_i=x_0+ih$. Define $s$ by $x=x_0+sh$, i.e.\ $s=\frac{x-x_0}{h}$.}
\why{$s$ counts how many steps of size $h$ you are from $x_0$: $s=0$ at $x_0$, $s=1$ at $x_1$, $s=2$ at $x_2$, and so on.}
\step{Each factor $x-x_i$ becomes a multiple of $h$:}
\[x-x_i=(x_0+sh)-(x_0+ih)=(s-i)\,h\]
\step{For integrals: $x=x_0+sh$ gives $\frac{dx}{ds}=h$, so $dx=h\,ds$, and the limits change from $x$-values to $s$-values:}
\[\int_{x_0}^{x_1}(\cdots)\,dx=\int_{0}^{1}(\cdots)\,h\,ds,\qquad\int_{x_0}^{x_2}(\cdots)\,dx=\int_{0}^{2}(\cdots)\,h\,ds\]
\result{$x-x_i=(s-i)h$,\quad $dx=h\,ds$,\quad $x_i\leftrightarrow s=i$.}
\end{deriv}

\begin{deriv}{T5}{Algebra identities that keep appearing}{\stwo}
\begin{align*}
p^2-q^2&=(p-q)(p+q) && \text{(cancellation fixes; secant for $\sqrt a$)}\\
p^3-q^3&=(p-q)(p^2+pq+q^2) && \text{(secant for $\sqrt[3]a$)}\\
1+r+r^2+\cdots+r^{k-1}&=\frac{1-r^k}{1-r}\quad(r\neq1) && \text{(largest floating point number)}\\
ax^2+2bxy&=a\Big(x+\frac bay\Big)^2-\frac{b^2}{a}y^2 && \text{(completing the square, positive definiteness)}
\end{align*}
\chk{expand the right-hand side of each line and you get the left-hand side.}
\end{deriv}

% ======================= WEEKS 1-2 =======================
\topic{Weeks 1--2: Floating point, cancellation and Taylor polynomials}

\begin{deriv}{D1.1}{Smallest positive number $S=\beta^m$}{\stwo\ the reasoning behind 2023 Q1b and 2025 Q1a}
\goal{show that the smallest positive number in $F(\beta,k,m,M)$ is $S=\beta^m$.}
\idea{a floating point number is (significand) $\times$ (base)$^{\text{exponent}}$. Make \emph{both} parts as small as the rules allow.}

\step{Every nonzero number in the system has the form}
\[\pm\,d_1.d_2d_3\cdots d_k\times\beta^{e},\qquad d_1\neq0,\qquad m\le e\le M\]
\why{$d_1\neq0$ is the ``normalised'' rule: it makes every number's representation unique.}

\step{Make the significand as small as possible. The first digit cannot be 0, so the smallest choice is $d_1=1$. Every other digit can be 0:}
\[d_1.d_2\cdots d_k=1.00\cdots0=1\]

\step{Make the exponent as small as possible: $e=m$, giving the factor $\beta^m$.}

\step{Multiply the two smallest parts:}
\[S=1\times\beta^{m}=\beta^{m}\]

\result{$S=\beta^m$. \quad Example: $F(2,5,-4,4)$ gives $S=2^{-4}=0.0625$.}
\why{any positive number smaller than $S$ would need an exponent below $m$, which isn't allowed. It \textbf{underflows} to 0.}
\end{deriv}

\begin{deriv}{D1.2}{Largest number $L=(\beta-\beta^{1-k})\,\beta^M$}{\stwo}
\goal{show that the largest number in $F(\beta,k,m,M)$ is $L=(\beta-\beta^{1-k})\beta^M$.}
\idea{make every digit as large as possible ($\beta-1$) and use the largest exponent $M$, then tidy up the significand with a geometric series.}

\step{The largest digit in base $\beta$ is $\beta-1$ (e.g.\ 9 in base 10, 1 in base 2). Use it for all $k$ digits, with exponent $M$:}
\[L=(\beta-1).(\beta-1)\cdots(\beta-1)\times\beta^{M}\qquad(k\text{ digits})\]

\step{Write the significand as a sum of place values. The digits sit in the places $\beta^0,\beta^{-1},\dots,\beta^{-(k-1)}$:}
\[(\beta-1)\Big(1+\beta^{-1}+\beta^{-2}+\cdots+\beta^{-(k-1)}\Big)\]

\step{The bracket is a geometric series with $k$ terms and ratio $r=\beta^{-1}$. Use $1+r+\cdots+r^{k-1}=\frac{1-r^k}{1-r}$ (Toolkit T5):}
\[1+\beta^{-1}+\cdots+\beta^{-(k-1)}=\frac{1-\beta^{-k}}{1-\beta^{-1}}=\frac{\beta\,(1-\beta^{-k})}{\beta-1}\]
\why{in the last step, multiply the top and bottom by $\beta$, so the bottom becomes $\beta-1$.}

\step{Multiply by $(\beta-1)$. The $(\beta-1)$ cancels:}
\[(\beta-1)\cdot\frac{\beta\,(1-\beta^{-k})}{\beta-1}=\beta(1-\beta^{-k})=\beta-\beta^{1-k}\]

\step{Attach the exponent:}
\[L=(\beta-\beta^{1-k})\,\beta^{M}\]

\result{$L=(\beta-\beta^{1-k})\beta^M$.}
\why{quick intuition: just as $9.99=10-0.01$, the largest significand is $\beta$ minus one unit in the last place, and that unit is $\beta^{1-k}$.}
\chk{$F(2,4,-2,3)$: $(2-2^{-3})\times2^3=1.875\times8=15$ \cmark. \quad IEEE double $F(2,53,-1022,1023)$: $(2-2^{-52})\times2^{1023}\approx1.80\times10^{308}$ \cmark}
\end{deriv}

\begin{deriv}{D1.3}{Why rounding error is at most $u=\tfrac12\beta^{1-k}$}{\sone\ understanding (the lectures state it without proof)}
\goal{show that $\dfrac{|x-fl(x)|}{|x|}\le\tfrac12\beta^{1-k}$ for any $x$ inside the range of $F$.}
\idea{rounding to the nearest float is wrong by at most \emph{half a gap}. Measure that gap relative to the size of $x$.}

\step{Write $x$ in normalised form with exponent $e$: $x=d_1.d_2d_3\cdots\times\beta^{e}$, where the significand is between $1$ and $\beta$.}

\step{The floats with exponent $e$ differ only in their last ($k$th) digit, so neighbouring floats are one unit in the last place apart. In the significand that unit is $\beta^{1-k}$, so the actual gap is}
\[\text{gap}=\beta^{1-k}\,\beta^{e}\]

\step{Rounding to the \emph{nearest} float means $x$ is at most half a gap away from $fl(x)$:}
\[|x-fl(x)|\le\tfrac12\,\beta^{1-k}\,\beta^{e}\]

\step{The significand is at least 1, so $|x|\ge1\times\beta^{e}=\beta^{e}$.}

\step{Divide the bound in Step 3 by $|x|$, using Step 4 (a smaller denominator can only make the right-hand side bigger):}
\[\frac{|x-fl(x)|}{|x|}\le\frac{\tfrac12\beta^{1-k}\beta^{e}}{\beta^{e}}=\tfrac12\beta^{1-k}=u\]

\result{$\dfrac{|x-fl(x)|}{|x|}\le u=\tfrac12\beta^{1-k}$, as long as $x$ doesn't overflow or underflow.}
\end{deriv}

\begin{deriv}{D2.1}{Removing catastrophic cancellation from the quadratic formula}{\sthree\ 2025 Q1b (3 marks)}
\goal{rewrite $x=\dfrac{-b+\sqrt{b^2-4ac}}{2a}$ so it has no catastrophic cancellation when $b^2\gg4ac$.}
\idea{multiply by the ``conjugate over itself'' (which equals 1) so the subtraction turns into an addition.}

\step{\mk{1} \textbf{Recognise the danger.} If $b^2\gg4ac$, then $b^2-4ac\approx b^2$, so $\sqrt{b^2-4ac}\approx|b|$. For $b>0$ the numerator is approximately $-b+b$: two nearly-equal numbers are subtracted, which is catastrophic cancellation.}

\step{\mk{1} \textbf{Multiply by the conjugate over itself.} Write $D=b^2-4ac$ to keep things short:}
\[x=\frac{-b+\sqrt D}{2a}\times\frac{-b-\sqrt D}{-b-\sqrt D}\]
\why{the second fraction equals 1, so the value of $x$ doesn't change, but the form does.}

\step{\textbf{Numerator: difference of squares.} With $p=-b$ and $q=\sqrt D$, $(p+q)(p-q)=p^2-q^2$:}
\[(-b+\sqrt D)(-b-\sqrt D)=(-b)^2-(\sqrt D)^2=b^2-(b^2-4ac)=4ac\]
\why{the $b^2$ terms cancel \emph{exactly in the algebra}. Exact algebraic cancellation is harmless; only cancellation of rounded numbers is dangerous.}

\step{\textbf{Denominator:} $2a(-b-\sqrt D)=-2a\,(b+\sqrt D)$.}

\step{\mk{1} \textbf{Simplify:}}
\[x=\frac{4ac}{-2a\,(b+\sqrt D)}=\frac{-2c}{b+\sqrt{b^2-4ac}}\]

\result{$x=\dfrac{-2c}{b+\sqrt{b^2-4ac}}$. The denominator \textbf{adds} two positive numbers, so there is no cancellation.}

\why{if instead $b<0$, it's the \emph{other} root $\frac{-b-\sqrt D}{2a}$ that is dangerous. The same trick (multiply by $\frac{-b+\sqrt D}{-b+\sqrt D}$) gives $\dfrac{-2c}{b-\sqrt D}$.}
\end{deriv}

\begin{deriv}{D2.2}{Other cancellation fixes}{\stwo\ lecture problems; likely exam variations}
\goal{rewrite each expression so it avoids cancellation in the case shown.}

\step{$1-\cos x$ for $x\approx0$ (since $\cos x\approx1$). Multiply by $\frac{1+\cos x}{1+\cos x}$ and use $1-\cos^2x=\sin^2x$:}
\[1-\cos x=\frac{(1-\cos x)(1+\cos x)}{1+\cos x}=\frac{1-\cos^2x}{1+\cos x}=\frac{\sin^2x}{1+\cos x}\]

\step{$\log(x+1)-\log x$ for large $x$ (two nearly-equal logs). Use the log rule $\log A-\log B=\log\frac AB$:}
\[\log(x+1)-\log x=\log\frac{x+1}{x}=\log\Big(1+\frac1x\Big)\]

\step{$\sqrt{x+1}-\sqrt x$ for large $x$. Use the conjugate:}
\[\sqrt{x+1}-\sqrt x=\frac{(x+1)-x}{\sqrt{x+1}+\sqrt x}=\frac{1}{\sqrt{x+1}+\sqrt x}\]

\step{$\sqrt{x^2+1}-x$ for large $x$. Use the conjugate:}
\[\sqrt{x^2+1}-x=\frac{(x^2+1)-x^2}{\sqrt{x^2+1}+x}=\frac{1}{\sqrt{x^2+1}+x}\]

\result{every fix turns ``subtract two nearly-equal numbers'' into ``add'' or ``divide''.}
\end{deriv}

\begin{deriv}{D2.3}{Where the Taylor polynomial's coefficients come from}{\stwo}
\goal{show that the polynomial matching $f,f',\dots,f^{(n)}$ at $x_0$ is $P_n(x)=\displaystyle\sum_{k=0}^n\frac{f^{(k)}(x_0)}{k!}(x-x_0)^k$.}
\idea{write a general polynomial in powers of $(x-x_0)$, then choose each coefficient so the derivatives match at $x_0$.}

\step{Write a general degree-$n$ polynomial in powers of $(x-x_0)$:}
\[P(x)=a_0+a_1(x-x_0)+a_2(x-x_0)^2+a_3(x-x_0)^3+\cdots+a_n(x-x_0)^n\]
\why{at $x=x_0$, every bracket is 0, so only one term survives each time. That makes the matching easy.}

\step{\textbf{Match the value.} Put $x=x_0$: every term except $a_0$ vanishes, so $P(x_0)=a_0$. We need $P(x_0)=f(x_0)$:}
\[a_0=f(x_0)\]

\step{\textbf{Match the first derivative.} Differentiate, then put $x=x_0$:}
\[P'(x)=a_1+2a_2(x-x_0)+3a_3(x-x_0)^2+\cdots\quad\Rightarrow\quad P'(x_0)=a_1=f'(x_0)\]

\step{\textbf{Match the second derivative.}}
\[P''(x)=2a_2+3\cdot2\,a_3(x-x_0)+\cdots\quad\Rightarrow\quad P''(x_0)=2a_2=f''(x_0)\quad\Rightarrow\quad a_2=\frac{f''(x_0)}{2!}\]

\step{\textbf{Match the third derivative.}}
\[P'''(x_0)=3\cdot2\cdot1\,a_3=3!\,a_3=f'''(x_0)\quad\Rightarrow\quad a_3=\frac{f'''(x_0)}{3!}\]

\step{\textbf{Spot the pattern.} Differentiating $(x-x_0)^k$ exactly $k$ times leaves $k!$, so in general}
\[a_k=\frac{f^{(k)}(x_0)}{k!}\]

\result{$P_n(x)=f(x_0)+f'(x_0)(x-x_0)+\dfrac{f''(x_0)}{2!}(x-x_0)^2+\cdots+\dfrac{f^{(n)}(x_0)}{n!}(x-x_0)^n$.}
\why{the remainder $R_n(x)=\frac{f^{(n+1)}(c)}{(n+1)!}(x-x_0)^{n+1}$ is \emph{Taylor's theorem}. MXB103 states it without proof. It looks like ``the next term'', except that the derivative is evaluated at an unknown $c$ between $x_0$ and $x$.}
\end{deriv}

\begin{deriv}{D2.4}{Taylor polynomial and error bound: the full exam answer}{\sthree\ Q2 every year (6 marks)}
\goal{(2025 Q2) (a) find $P_2$ for $f(x)=\cos x$ centred on $x_0=0$; \ (b) bound the error for $0\le x\le\pi/3$.}
\idea{(a) differentiate, evaluate at $x_0$, substitute. \ (b) write the remainder, then replace each unknown by its worst case.}

\textbf{Part (a)}

\step{\mk{1} Write the form of the Taylor polynomial:}
\[P_2(x)=f(x_0)+f'(x_0)(x-x_0)+\frac{f''(x_0)}{2}(x-x_0)^2\]

\step{\mk{1} Differentiate and evaluate at $x_0=0$:}
\[f(x)=\cos x\Rightarrow f(0)=1,\qquad f'(x)=-\sin x\Rightarrow f'(0)=0,\qquad f''(x)=-\cos x\Rightarrow f''(0)=-1\]

\step{\mk{1} Substitute:}
\[P_2(x)=1+0\cdot x+\frac{-1}{2}x^2=1-\frac{x^2}{2}\]

\textbf{Part (b)}

\step{\mk{1} Write the remainder, and say where $c$ is. The next derivative is $f'''(x)=\sin x$:}
\[R_2(x)=\frac{f'''(c)}{3!}(x-0)^3=\frac{\sin c}{6}\,x^3,\qquad c\text{ between }0\text{ and }x\]
\why{$x$ is between $0$ and $\pi/3$, so $c$ is too.}

\step{\mk{1} Bound the derivative factor. On $[0,\pi/3]$, $\sin$ is increasing, so its largest value is at the right end:}
\[|\sin c|\le\sin\tfrac\pi3=\tfrac{\sqrt3}{2}\]

\step{Bound the distance factor. The farthest point of $[0,\pi/3]$ from $x_0=0$ is $\pi/3$:}
\[|x|^3\le\Big(\frac\pi3\Big)^3\]

\step{\mk{1} Multiply:}
\[|f(x)-P_2(x)|=|R_2(x)|\le\frac16\cdot\frac{\sqrt3}{2}\cdot\Big(\frac\pi3\Big)^3=0.1658\]

\result{(a) $P_2(x)=1-\dfrac{x^2}2$. \quad (b) Absolute error $\le0.1658$ for all $0\le x\le\pi/3$.}

\trap{when $x_0\neq0$, the distance factor is the farthest distance \emph{from $x_0$}, not the right endpoint. In 2023 Q2 ($x_0=\pi/6$ on $[0,\pi/3]$) it is $\max(|0-\frac\pi6|,|\frac\pi3-\frac\pi6|)=\frac\pi6$, giving the bound $\frac12\big(\frac\pi6\big)^2=\frac{\pi^2}{72}$.}
\end{deriv}

% ======================= WEEKS 3-4 =======================
\topic{Weeks 3--4: ODE methods}

\textbf{Setting for every derivation in this section.} We solve $\dfrac{dy}{dt}=f(t,y)$ on $a\le t\le b$ with $y(a)=\alpha$. We use the grid $h=\frac{b-a}{n}$, $t_i=a+ih$, so $t_{i+1}-t_i=h$. We write $y_i=y(t_i)$ for the exact (unknown) values and $w_i\approx y_i$ for the computed ones.

\textbf{The same four moves appear in every ODE derivation:}
\begin{center}
\begin{tabular}{cl}
\toprule
1 & Taylor-expand $y(t)$ about $t_i$\\
2 & Evaluate at $t_{i+1}$, so every $(t-t_i)$ becomes $h$\\
3 & Use the ODE to replace derivatives of $y$ by $f$\\
4 & Drop the $O(h^{\cdot})$ term and replace $y_i$ by $w_i$\\
\bottomrule
\end{tabular}
\end{center}

\begin{deriv}{D3.1}{Euler's method}{\stwo\ the template for every ODE derivation}
\goal{derive $w_{i+1}=w_i+h\,f(t_i,w_i)$.}
\idea{a degree-1 Taylor polynomial is the tangent line. Follow the tangent for one step.}

\step{\textbf{Move 1.} Taylor polynomial of degree 1 for $y(t)$ about $t_i$, with remainder:}
\[y(t)=y(t_i)+y'(t_i)(t-t_i)+\frac{y''(c)}{2}(t-t_i)^2\]
\why{we know $y$ at $t_i$ and want it at the next time. Taylor predicts nearby values from information at one point.}

\step{\textbf{Move 2.} Evaluate at $t=t_{i+1}$. Since $t_{i+1}-t_i=h$:}
\[y_{i+1}=y_i+h\,y'(t_i)+\frac{y''(c)}{2}h^2=y_i+h\,y'(t_i)+O(h^2)\]
\why{$c$ is unknown but $y''$ is bounded, so the last term is at most a constant times $h^2$ (Toolkit T2).}

\step{\textbf{Move 3.} The ODE says $y'(t)=f(t,y(t))$. At $t=t_i$ this gives $y'(t_i)=f(t_i,y_i)$:}
\[y_{i+1}=y_i+h\,f(t_i,y_i)+O(h^2)\]

\step{\textbf{Move 4.} Drop the $O(h^2)$ term, and replace the exact $y_i$ by the approximations $w_i$:}
\[w_{i+1}=w_i+h\,f(t_i,w_i),\qquad w_0=\alpha\]

\result{Euler's method: $w_{i+1}=w_i+h\,f(t_i,w_i)$.}
\why{the term we dropped, $O(h^2)$, is exactly the error made in one step: the \textbf{local error}.}
\end{deriv}

\begin{deriv}{D3.2}{Local error $O(h^{p+1})$ gives global error $O(h^p)$}{\stwo\ explains ``Euler is first order''}
\goal{show why a method whose single step has error $O(h^{p+1})$ has total error $O(h^p)$.}
\idea{the total error is roughly (error per step) $\times$ (number of steps), and the number of steps is $\frac{b-a}{h}$.}

\step{Each step adds an error of at most $C\,h^{p+1}$, for some constant $C$ (the local error).}
\step{Getting from $a$ to $b$ takes $n=\dfrac{b-a}{h}$ steps.}
\step{Add up all the steps' errors (roughly, ignoring how earlier errors are carried forward):}
\[\text{global error}\approx n\times C\,h^{p+1}=\frac{b-a}{h}\times C\,h^{p+1}=C(b-a)\,h^{p}=O(h^p)\]
\why{one power of $h$ is lost because a smaller $h$ means \emph{more} steps. A careful proof also tracks how earlier errors propagate, but the conclusion is the same.}

\result{local $O(h^{p+1})$ $\Rightarrow$ global $O(h^p)$. Euler: local $O(h^2)$, so global $O(h)$ (\textbf{first order}).}
\end{deriv}

\begin{deriv}{D4.1}{The total derivative $f'(t,y)$}{\sthree\ every Taylor 2 question: 2025 gave 0 marks if this is wrong}
\goal{show $f'(t,y)=\dfrac{\partial f}{\partial t}+\dfrac{\partial f}{\partial y}\cdot f(t,y)$, and that $y''=f'(t,y)$.}
\idea{along the solution, $y$ is itself a function of $t$, so $t$ affects $f$ in two ways: directly, and through $y$.}

\step{Along the solution curve, $y=y(t)$. So $f(t,y(t))$ is really a function of $t$ alone, but $t$ enters it twice: in the first slot directly, and in the second slot through $y(t)$.}

\step{The chain rule for a function of two variables adds up both effects:}
\[\frac{d}{dt}f\big(t,y(t)\big)=\frac{\partial f}{\partial t}\cdot\frac{dt}{dt}+\frac{\partial f}{\partial y}\cdot\frac{dy}{dt}\]

\step{Simplify: $\frac{dt}{dt}=1$, and $\frac{dy}{dt}=y'=f(t,y)$ by the ODE:}
\[f'(t,y)=\frac{\partial f}{\partial t}+\frac{\partial f}{\partial y}\cdot f(t,y)\]

\step{Since $y'=f(t,y)$, differentiating once more gives $y''=\frac{d}{dt}f(t,y)=f'(t,y)$.}

\result{$f'(t,y)=f_t+f_y\cdot f$, \quad and \quad $y''=f'(t,y)$.}

\textbf{Worked (lecture): $f(t,y)=3t-\dfrac yt$.}
\begin{align*}
\frac{\partial f}{\partial t}&=3+\frac{y}{t^2} &&\text{(treat $y$ as a constant: $\tfrac{d}{dt}(-yt^{-1})=yt^{-2}$)}\\
\frac{\partial f}{\partial y}&=-\frac1t &&\text{(treat $t$ as a constant)}\\
f'(t,y)&=3+\frac{y}{t^2}+\Big(-\frac1t\Big)\Big(3t-\frac yt\Big) &&\text{($f_t+f_y\cdot f$)}\\
&=3+\frac{y}{t^2}-3+\frac{y}{t^2}=\frac{2y}{t^2}
\end{align*}
\trap{writing $f'=3+\frac{y}{t^2}$, i.e.\ forgetting the $f_y\cdot f$ part. In MATLAB this is the missing \texttt{*f(t,y)} bug (2023 Q4, 2024 Q4).}
\end{deriv}

\begin{deriv}{D4.2}{Second-order Taylor method}{\sthree\ 2023 Q5 (4 marks)}
\goal{derive $w_{i+1}=w_i+h\,f(t_i,w_i)+\dfrac{h^2}{2}f'(t_i,w_i)$.}
\idea{the same four moves as Euler, but keep one more Taylor term.}

\step{\mk{1} \textbf{Move 1.} Taylor polynomial of degree 2 for $y(t)$ about $t_i$:}
\[y(t)=y(t_i)+y'(t_i)(t-t_i)+\frac{y''(t_i)}{2}(t-t_i)^2+O\big((t-t_i)^3\big)\]

\step{\mk{1} \textbf{Move 2.} Evaluate at $t=t_{i+1}$, so $t_{i+1}-t_i=h$:}
\[y_{i+1}=y_i+h\,y'(t_i)+\frac{h^2}{2}y''(t_i)+O(h^3)\]

\step{\mk{1} \textbf{Move 3.} Replace the derivatives of $y$ using the ODE and D4.1: $y'=f(t,y)$ and $y''=f'(t,y)$:}
\[y_{i+1}=y_i+h\,f(t_i,y_i)+\frac{h^2}{2}f'(t_i,y_i)+O(h^3)\]

\step{\mk{1} \textbf{Move 4.} Drop $O(h^3)$, and replace $y_i$ with $w_i$:}
\[w_{i+1}=w_i+h\,f(t_i,w_i)+\frac{h^2}{2}f'(t_i,w_i)\]

\result{second-order Taylor method. Local error $O(h^3)$, so global $O(h^2)$ (second order).}
\end{deriv}

\begin{deriv}{D4.3}{Modified Euler method}{\sthree\ 2024 Q5 and 2025 Q5 (4 marks each)}
\goal{derive $w_{i+1}=w_i+\dfrac h2\Big[f(t_i,w_i)+f\big(t_{i+1},\,w_i+hf(t_i,w_i)\big)\Big]$.}
\idea{start from Taylor 2, but get rid of the awkward $f'$ by approximating it with a difference quotient (a slope). Then patch the one remaining problem using Euler.}

\step{\mk{0.5} \textbf{Taylor polynomial of degree 2} for $y(t)$ about $t_i$:}
\[y(t)=y(t_i)+y'(t_i)(t-t_i)+\frac{y''(t_i)}{2}(t-t_i)^2+O\big((t-t_i)^3\big)\]

\step{\mk{0.5} \textbf{Evaluate at $t_{i+1}$}, so $t_{i+1}-t_i=h$:}
\[y_{i+1}=y_i+h\,y'(t_i)+\frac{h^2}{2}y''(t_i)+O(h^3)\]

\step{\mk{0.5} \textbf{Replace the derivatives of $y$ by $f$} ($y'=f$, $y''=f'$):}
\[y_{i+1}=y_i+h\,f(t_i,y_i)+\frac{h^2}{2}f'(t_i,y_i)+O(h^3)\]
\why{up to here this is just Taylor 2. Next we remove $f'$, because computing it needs the chain rule.}

\step{\mk{0.5} \textbf{Approximate $f'$ by a forward difference}, the slope of $f$ between $t_i$ and $t_{i+1}$ (Week 10, D10.1):}
\[f'(t_i,y_i)=\frac{f(t_{i+1},y_{i+1})-f(t_i,y_i)}{h}+O(h)\]

\step{\mk{1} \textbf{Substitute and simplify.} To keep it readable, write $f_i=f(t_i,y_i)$ and $f_{i+1}=f(t_{i+1},y_{i+1})$:}
\begin{align*}
y_{i+1}&=y_i+h f_i+\frac{h^2}{2}\Big[\frac{f_{i+1}-f_i}{h}+O(h)\Big]+O(h^3)\\
&=y_i+h f_i+\frac h2\,(f_{i+1}-f_i)+O(h^3) &&\text{($\tfrac{h^2}{2}\cdot\tfrac1h=\tfrac h2$, and $\tfrac{h^2}{2}\cdot O(h)=O(h^3)$)}\\
&=y_i+\frac h2 f_i+\frac h2 f_{i+1}+O(h^3) &&\text{($h f_i-\tfrac h2 f_i=\tfrac h2 f_i$)}\\
&=y_i+\frac h2\big[f(t_i,y_i)+f(t_{i+1},y_{i+1})\big]+O(h^3)
\end{align*}

\step{\textbf{Drop $O(h^3)$ and replace $y$ by $w$:}}
\[w_{i+1}=w_i+\frac h2\Big[f(t_i,w_i)+f(t_{i+1},\,w_{i+1})\Big]\]

\step{\mk{1} \textbf{Patch the problem: $w_{i+1}$ appears on both sides.} We are trying to \emph{compute} $w_{i+1}$, so it can't be on the right-hand side. Inside $f$ on the right, replace it with Euler's prediction $w_{i+1}\approx w_i+h\,f(t_i,w_i)$:}
\[w_{i+1}=w_i+\frac h2\Big[f(t_i,w_i)+f\big(t_{i+1},\,w_i+h\,f(t_i,w_i)\big)\Big]\]
\why{Euler's error inside $f$ is $O(h^2)$. It's multiplied by $\frac h2$, so it only adds $O(h^3)$, which is the same size as the error we already dropped. The method stays second order.}

\step{\textbf{(Optional) the $k$-form.} Name the two increments $k_1=h\,f(t_i,w_i)$ and $k_2=h\,f(t_i+h,\,w_i+k_1)$. Then $\frac h2[\,\cdot\,]=\frac12(k_1+k_2)$:}
\[w_{i+1}=w_i+\tfrac12(k_1+k_2)\]

\result{Modified Euler: $w_{i+1}=w_i+\frac h2\big[f(t_i,w_i)+f(t_{i+1},w_i+hf(t_i,w_i))\big]$. Global $O(h^2)$, with no $f'$ needed.}
\end{deriv}

\begin{deriv}{D4.4}{RK4: what you need to know (no derivation required)}{\sone}
\goal{know the RK4 formulas, and see why they make sense.}
\step{The lectures state RK4 without derivation (it takes ``laborious algebra''). You need the formulas and the order (global $O(h^4)$) only:}
\begin{align*}
k_1&=hf(t_i,w_i), & k_2&=hf\big(t_i+\tfrac h2,\,w_i+\tfrac{k_1}{2}\big),\\
k_3&=hf\big(t_i+\tfrac h2,\,w_i+\tfrac{k_2}{2}\big), & k_4&=hf(t_i+h,\,w_i+k_3)
\end{align*}
\[w_{i+1}=w_i+\tfrac16(k_1+2k_2+2k_3+k_4)\]

\step{\textbf{Sanity check: if $f$ depends only on $t$}, then $k_2=k_3=h f(t_i+\frac h2)$, and}
\[w_{i+1}-w_i=\frac h6\Big[f(t_i)+4f\big(t_i+\tfrac h2\big)+f(t_{i+1})\Big]\]
\why{that is exactly Simpson's rule (D11.3) on $[t_i,t_{i+1}]$, using the strip half-width $\frac h2$: $\frac{h/2}{3}[f+4f+f]$. So RK4 really is ``Simpson's rule for ODEs'', which is why the weights are $1,2,2,1$.}
\result{RK4 is a weighted average of four slope samples, and reduces to Simpson's rule when $f=f(t)$.}
\end{deriv}

% ======================= WEEKS 5-7 =======================
\topic{Weeks 5--7: Interpolation}

\textbf{Setting.} We have data points $(x_0,y_0),(x_1,y_1),\dots,(x_n,y_n)$ with \emph{distinct} $x_i$. We want the polynomial $P_n$ of degree $\le n$ with $P_n(x_i)=y_i$ for every $i$. The key theorem: \textbf{such a polynomial exists and is unique}. Every ``form'' below is the same polynomial, written differently.

\begin{deriv}{D5.1}{From two points to the Lagrange form}{\stwo\ lecture derivation; shows where Lagrange comes from}
\goal{find the line through $(x_0,y_0)$ and $(x_1,y_1)$, and rearrange it into $P_1(x)=\dfrac{x-x_1}{x_0-x_1}y_0+\dfrac{x-x_0}{x_1-x_0}y_1$.}
\idea{solve for the line directly, then regroup the result by $y_0$ and $y_1$.}

\step{Look for $P_1(x)=a_0+a_1x$. It must pass through both points:}
\[a_0+a_1x_0=y_0\quad(1)\qquad\qquad a_0+a_1x_1=y_1\quad(2)\]

\step{Subtract (1) from (2) to eliminate $a_0$:}
\[a_1(x_1-x_0)=y_1-y_0\quad\Rightarrow\quad a_1=\frac{y_1-y_0}{x_1-x_0}\quad\text{(the slope)}\]

\step{Substitute back into (1) and solve for $a_0$:}
\[a_0=y_0-a_1x_0=\frac{y_0(x_1-x_0)-(y_1-y_0)x_0}{x_1-x_0}=\frac{x_1y_0-x_0y_1}{x_1-x_0}\]

\step{Put both into $P_1$ over a common denominator, then collect the $y_0$ terms and the $y_1$ terms:}
\[P_1(x)=\frac{x_1y_0-x_0y_1+(y_1-y_0)x}{x_1-x_0}=\frac{y_0\,(x_1-x)+y_1\,(x-x_0)}{x_1-x_0}\]

\step{Split into two fractions. For the $y_0$ term, flip the signs of the top and bottom ($\frac{x_1-x}{x_1-x_0}=\frac{x-x_1}{x_0-x_1}$):}
\[P_1(x)=\underbrace{\frac{x-x_1}{x_0-x_1}}_{L_{1,0}(x)}y_0+\underbrace{\frac{x-x_0}{x_1-x_0}}_{L_{1,1}(x)}y_1\]

\result{$P_1(x)=L_{1,0}(x)\,y_0+L_{1,1}(x)\,y_1$. Note that $L_{1,0}$ equals 1 at $x_0$ and 0 at $x_1$, and $L_{1,1}$ does the opposite. That pattern is what generalises.}
\end{deriv}

\begin{deriv}{D5.2}{Why the Lagrange form interpolates}{\stwo}
\goal{show that $P_n(x)=\sum_{i=0}^nL_{n,i}(x)\,y_i$ with $L_{n,i}(x)=\prod_{j\neq i}\frac{x-x_j}{x_i-x_j}$ passes through every data point.}
\idea{each $L_{n,i}$ is designed to be 1 at its own node and 0 at all the others.}

\step{\textbf{$L_{n,i}$ is zero at the other nodes.} Its numerator $\prod_{j\ne i}(x-x_j)$ contains the factor $(x-x_k)$ for every $k\neq i$. So at $x=x_k$ that factor is 0:}
\[L_{n,i}(x_k)=0\qquad(k\neq i)\]

\step{\textbf{$L_{n,i}$ is one at its own node.} At $x=x_i$ the numerator becomes $\prod_{j\ne i}(x_i-x_j)$, which is exactly the denominator:}
\[L_{n,i}(x_i)=\frac{\prod_{j\ne i}(x_i-x_j)}{\prod_{j\ne i}(x_i-x_j)}=1\]
\why{this is why the nodes must be distinct: if $x_i=x_j$, the denominator would be 0.}

\step{\textbf{Evaluate $P_n$ at a node $x_k$.} Every term with $i\neq k$ vanishes, and the $i=k$ term has $L_{n,k}(x_k)=1$:}
\[P_n(x_k)=\sum_iL_{n,i}(x_k)\,y_i=L_{n,k}(x_k)\,y_k=y_k\]

\step{\textbf{Degree.} Each $L_{n,i}$ is a product of $n$ linear factors, so it has degree $n$. A sum of these has degree at most $n$.}

\result{$P_n$ has degree $\le n$ and $P_n(x_k)=y_k$ for every $k$, so it is \emph{the} interpolating polynomial.}

\textbf{Written out for three points:}
\[P_2(x)=\frac{(x-x_1)(x-x_2)}{(x_0-x_1)(x_0-x_2)}y_0+\frac{(x-x_0)(x-x_2)}{(x_1-x_0)(x_1-x_2)}y_1+\frac{(x-x_0)(x-x_1)}{(x_2-x_0)(x_2-x_1)}y_2\]
\end{deriv}

\begin{deriv}{D5.3}{Why the Lagrange basis polynomials add up to 1}{\sone\ the arithmetic check}
\goal{show $\sum_iL_{n,i}(x)=1$ for every $x$.}
\step{Interpolate the constant function $f(x)=1$. Then every $y_i=1$, so the Lagrange form gives $P_n(x)=\sum_iL_{n,i}(x)\cdot1$.}
\step{The constant polynomial $1$ also passes through all these points, and it has degree $0\le n$.}
\step{By uniqueness, these are the same polynomial: $\sum_iL_{n,i}(x)=1$.}
\result{$\sum_iL_{n,i}(x)=1$. If your computed $L$ values don't add to 1, one of them is wrong.}
\end{deriv}

\begin{deriv}{D5.4}{Using the interpolation remainder (two cases)}{\stwo\ 2024 Q6b (3 marks)}
\goal{bound $|f(x)-P_n(x)|$ using $R_n(x)=\dfrac{f^{(n+1)}(c)}{(n+1)!}(x-x_0)(x-x_1)\cdots(x-x_n)$, which is stated without proof.}
\step{Write $R_n$, with $c$ in the interval containing the nodes and $x$.}
\step{Bound $|f^{(n+1)}(c)|$ by its \textbf{maximum} over that interval ($c$ is unknown, so use the worst case).}
\step{Handle the product:}
\begin{itemize}[leftmargin=18pt,itemsep=1pt]
\item error at \textbf{one point} $\bar x$: just \emph{evaluate} $|(\bar x-x_0)\cdots(\bar x-x_n)|$;
\item error for \textbf{all $x$ in an interval}: use the \emph{maximum} of the product (usually given in the exam).
\end{itemize}
\step{Multiply and divide by $(n+1)!$.}

\textbf{Worked (2024 Q6b):} $f=\cos x$, nodes $0,\frac\pi6,\frac\pi2$, all $x\in[0,\frac\pi4]$, given $|x(x-\frac\pi6)(x-\frac\pi2)|<0.17$:
\[|R_2(x)|=\Big|\frac{\sin c}{3!}\Big|\,\Big|x\Big(x-\frac\pi6\Big)\Big(x-\frac\pi2\Big)\Big|<\frac{1}{6}\times0.17=0.0283\]
\why{$f'''(x)=\sin x$, and $|\sin c|\le1$.}
\result{error $\le\dfrac{\max|f^{(n+1)}|}{(n+1)!}\times|\text{product}|$ (evaluated at one point, or its max over the interval).}
\end{deriv}

\begin{deriv}{D6.1}{Newton divided difference coefficients}{\stwo\ lecture derivation}
\goal{show that the coefficients of $P_n(x)=a_0+a_1(x-x_0)+a_2(x-x_0)(x-x_1)+\cdots$ are the divided differences $a_k=f[x_0,\dots,x_k]$, by checking $a_0$, $a_1$ and $a_2$.}
\idea{impose $P(x_0)=y_0$, $P(x_1)=y_1$, $P(x_2)=y_2$ one at a time. Each new condition brings in only one new coefficient.}

\step{\textbf{At $x_0$:} every term except $a_0$ contains $(x-x_0)$, so they vanish:}
\[a_0=y_0=f[x_0]\]

\step{\textbf{At $x_1$:} terms from $a_2$ onwards contain $(x-x_1)$, so they vanish:}
\[a_0+a_1(x_1-x_0)=y_1\quad\Rightarrow\quad a_1=\frac{y_1-y_0}{x_1-x_0}=f[x_0,x_1]\]

\step{\textbf{At $x_2$:} terms from $a_3$ onwards vanish. Solve for $a_2$:}
\[a_2=\frac{y_2-y_0-a_1(x_2-x_0)}{(x_2-x_0)(x_2-x_1)}\]

\step{Simplify the numerator. Split $y_2-y_0=(y_2-y_1)+(y_1-y_0)$ and use $a_1=\frac{y_1-y_0}{x_1-x_0}$:}
\begin{align*}
y_2-y_0-a_1(x_2-x_0)&=(y_2-y_1)+(y_1-y_0)\Big[1-\frac{x_2-x_0}{x_1-x_0}\Big]\\
&=(y_2-y_1)+(y_1-y_0)\cdot\frac{x_1-x_2}{x_1-x_0}=(y_2-y_1)-(y_1-y_0)\frac{x_2-x_1}{x_1-x_0}
\end{align*}

\step{Divide by $(x_2-x_0)(x_2-x_1)$. Dividing the first piece by $(x_2-x_1)$ gives a slope, and the $(x_2-x_1)$ cancels in the second:}
\[a_2=\frac{\dfrac{y_2-y_1}{x_2-x_1}-\dfrac{y_1-y_0}{x_1-x_0}}{x_2-x_0}=\frac{f[x_1,x_2]-f[x_0,x_1]}{x_2-x_0}=f[x_0,x_1,x_2]\]

\result{$a_0=f[x_0]$, $a_1=f[x_0,x_1]$, $a_2=f[x_0,x_1,x_2]$, and in general $a_k=f[x_0,\dots,x_k]$, the top diagonal of the table.}
\end{deriv}

\begin{deriv}{D6.2}{Expanding $f[x_0,x_1,x_2,x_3]$}{\sthree\ 2023 Q8 and 2025 Q8 (4 marks, 1 per line)}
\goal{express $f[x_0,x_1,x_2,x_3]$ using only the $x_i$ and $y_i$.}
\idea{apply the definition (drop the first point) $-$ (drop the last point), all over (last $x$ $-$ first $x$), again and again until only $f[x_j]=y_j$ remain.}

\step{\mk{1} Apply the definition to the whole thing:}
\[f[x_0,x_1,x_2,x_3]=\frac{f[x_1,x_2,x_3]-f[x_0,x_1,x_2]}{x_3-x_0}\]

\step{\mk{1} Apply it to each second divided difference:}
\[=\frac{\dfrac{f[x_2,x_3]-f[x_1,x_2]}{x_3-x_1}-\dfrac{f[x_1,x_2]-f[x_0,x_1]}{x_2-x_0}}{x_3-x_0}\]

\step{\mk{1} Apply it to each first divided difference, $f[x_j,x_{j+1}]=\dfrac{f[x_{j+1}]-f[x_j]}{x_{j+1}-x_j}$:}
\[=\frac{\dfrac{\frac{f[x_3]-f[x_2]}{x_3-x_2}-\frac{f[x_2]-f[x_1]}{x_2-x_1}}{x_3-x_1}-\dfrac{\frac{f[x_2]-f[x_1]}{x_2-x_1}-\frac{f[x_1]-f[x_0]}{x_1-x_0}}{x_2-x_0}}{x_3-x_0}\]

\step{\mk{1} Replace each $f[x_j]$ by $y_j$:}
\[f[x_0,x_1,x_2,x_3]=\frac{\dfrac{\frac{y_3-y_2}{x_3-x_2}-\frac{y_2-y_1}{x_2-x_1}}{x_3-x_1}-\dfrac{\frac{y_2-y_1}{x_2-x_1}-\frac{y_1-y_0}{x_1-x_0}}{x_2-x_0}}{x_3-x_0}\]

\result{the 4-line expansion above. Write every level on its own line: each one is a mark.}
\trap{dividing by the neighbours' gap (e.g.\ $x_3-x_2$) at the second level. A second divided difference over $x_1,x_2,x_3$ divides by the \emph{outer} gap $x_3-x_1$.}
\end{deriv}

\begin{deriv}{D7.1}{Second and third forward differences}{\stwo}
\goal{express $\Delta^2y_i$ and $\Delta^3y_i$ in terms of the $y$'s.}
\idea{$\Delta$ means ``next minus current''. Apply it repeatedly.}
\step{Definition: $\Delta y_i=y_{i+1}-y_i$.}
\step{Second difference = $\Delta$ applied to $\Delta y_i$:}
\[\Delta^2y_i=\Delta y_{i+1}-\Delta y_i=(y_{i+2}-y_{i+1})-(y_{i+1}-y_i)=y_{i+2}-2y_{i+1}+y_i\]
\step{Third difference = $\Delta$ applied to $\Delta^2y_i$:}
\[\Delta^3y_i=\Delta^2y_{i+1}-\Delta^2y_i=(y_{i+3}-2y_{i+2}+y_{i+1})-(y_{i+2}-2y_{i+1}+y_i)=y_{i+3}-3y_{i+2}+3y_{i+1}-y_i\]
\result{the coefficients are Pascal's triangle with alternating signs: $1,-2,1$ and $1,-3,3,-1$.}
\end{deriv}

\begin{deriv}{D7.2}{Divided differences equal scaled forward differences}{\sthree\ 2024 Q8 (4 marks)}
\goal{for equally spaced nodes (spacing $h$), show $f[x_0,x_1,x_2]=\dfrac{\Delta^2y_0}{2!\,h^2}$.}
\idea{write out the divided difference, then use equal spacing: every gap between neighbours is $h$, and the outer gap is $2h$.}

\step{\textbf{Warm-up, $k=1$.} $f[x_0,x_1]=\dfrac{y_1-y_0}{x_1-x_0}=\dfrac{\Delta y_0}{h}=\dfrac{\Delta y_0}{1!\,h^1}$ \cmark}

\step{\mk{1} Apply the definition:}
\[f[x_0,x_1,x_2]=\frac{f[x_1,x_2]-f[x_0,x_1]}{x_2-x_0}\]

\step{\mk{1} Write out the first divided differences:}
\[=\frac{\dfrac{y_2-y_1}{x_2-x_1}-\dfrac{y_1-y_0}{x_1-x_0}}{x_2-x_0}\]

\step{Use equal spacing: $x_2-x_1=x_1-x_0=h$ and $x_2-x_0=2h$:}
\[=\frac{\dfrac{y_2-y_1}{h}-\dfrac{y_1-y_0}{h}}{2h}\]

\step{\mk{2} Simplify. Multiply the top and bottom by $h$, then recognise $\Delta^2y_0$ (D7.1):}
\[=\frac{(y_2-y_1)-(y_1-y_0)}{2h^2}=\frac{y_2-2y_1+y_0}{2h^2}=\frac{\Delta^2y_0}{2!\,h^2}\]

\result{$f[x_0,x_1,x_2]=\dfrac{\Delta^2y_0}{2!\,h^2}$. In general $f[x_0,\dots,x_k]=\dfrac{\Delta^ky_0}{k!\,h^k}$.}
\end{deriv}

\begin{deriv}{D7.3}{The Newton forward difference form}{\stwo}
\goal{for equally spaced nodes, turn the divided-difference form into $P_n(x)=\displaystyle\sum_{k=0}^n\binom sk\Delta^ky_0$ with $s=\frac{x-x_0}{h}$.}
\idea{substitute $x=x_0+sh$ (Toolkit T4). Every power of $h$ cancels.}

\step{The divided difference form (D6.1): term $k$ is $f[x_0,\dots,x_k]\,(x-x_0)(x-x_1)\cdots(x-x_{k-1})$.}

\step{With $x=x_0+sh$, each factor is $x-x_i=(s-i)h$, so the product of the first $k$ factors is}
\[(x-x_0)(x-x_1)\cdots(x-x_{k-1})=h^k\,s(s-1)(s-2)\cdots(s-k+1)\]

\step{The coefficient, from D7.2: $f[x_0,\dots,x_k]=\dfrac{\Delta^ky_0}{k!\,h^k}$.}

\step{Multiply them. The $h^k$ cancels:}
\[\frac{\Delta^ky_0}{k!\,h^k}\cdot h^k\,s(s-1)\cdots(s-k+1)=\frac{s(s-1)\cdots(s-k+1)}{k!}\,\Delta^ky_0=\binom sk\Delta^ky_0\]

\result{$P_n(x)=y_0+s\,\Delta y_0+\dfrac{s(s-1)}{2!}\Delta^2y_0+\dfrac{s(s-1)(s-2)}{3!}\Delta^3y_0+\cdots$, \quad $s=\dfrac{x-x_0}{h}$.}
\end{deriv}

% ======================= WEEKS 8-9 =======================
\topic{Weeks 8--9: Root finding}

\textbf{Setting.} We want $x$ with $f(x)=0$. Any equation can be put in this form by moving everything to one side.

\begin{deriv}{D8.1}{Showing $[a,b]$ is a bracketing}{\sthree\ Q9a every year (2 marks)}
\goal{show that $f$ has a root in $[a,b]$.}
\idea{if a continuous function is negative at one end and positive at the other, it must cross zero in between (Intermediate Value Theorem).}
\step{\mk{0.5 + 0.5} Evaluate $f$ at both ends.}
\step{\mk{1} State: they have \textbf{opposite signs}, and $f$ is \textbf{continuous} on $[a,b]$, so by the IVT $f$ takes every value between $f(a)$ and $f(b)$, including 0. Hence $[a,b]$ is a bracketing.}

\textbf{Worked (2025 Q9a):} $f(x)=x^3-2x-7$ on $[1,3]$.
\[f(1)=1-2-7=-8<0,\qquad f(3)=27-6-7=14>0\]
Opposite signs, and $f$ is a polynomial (so continuous). Hence $[1,3]$ is a bracketing for $f$.
\result{``$f(a)$ and $f(b)$ have opposite signs and $f$ is continuous, so $[a,b]$ brackets a root.''}
\end{deriv}

\begin{deriv}{D8.2}{Bisection error bound}{\stwo}
\goal{show that after $n$ iterations the midpoint $p_n$ satisfies $|p_n-r|\le\dfrac{b-a}{2^n}$, where $r$ is the root.}
\idea{every iteration halves the bracket, and the root is always inside it.}
\step{The starting bracket has length $b-a$, and each iteration halves it. So the bracket used in iteration $n$ has length $\dfrac{b-a}{2^{n-1}}$.}
\step{$p_n$ is the midpoint of that bracket, and the root $r$ is somewhere inside it. A point inside an interval is at most half its length from the midpoint:}
\[|p_n-r|\le\frac12\cdot\frac{b-a}{2^{n-1}}=\frac{b-a}{2^n}\]
\step{\textbf{How many iterations for a tolerance?} We need $\frac{b-a}{2^n}\le\text{tol}$, i.e.\ $2^n\ge\frac{b-a}{\text{tol}}$, so $n\ge\log_2\!\frac{b-a}{\text{tol}}$.}
\result{$|p_n-r|\le\dfrac{b-a}{2^n}$. \ E.g.\ on $[1,3]$ with tol $10^{-3}$: $2^n\ge2000$, so $n=11$.}
\end{deriv}

\begin{deriv}{D8.3}{Rearranging $f(x)=0$ into $x=g(x)$}{\stwo\ lecture example}
\goal{find two fixed point forms of $0.9\cos x=\sqrt x$ on $[0,1]$.}
\idea{any algebra that isolates an $x$ on one side gives a valid $g$. Different choices behave differently.}
\step{\textbf{Rearrangement 1.} Divide by 0.9, then apply $\cos^{-1}$:}
\[\cos x=\frac{\sqrt x}{0.9}\quad\Rightarrow\quad x=\cos^{-1}\!\Big(\frac{\sqrt x}{0.9}\Big)=g_1(x)\]
\step{\textbf{Rearrangement 2.} Square both sides (safe here, since both sides are $\ge0$ on $[0,1]$):}
\[0.81\cos^2x=x\quad\Rightarrow\quad x=0.81\cos^2x=g_2(x)\]
\result{same equation, two $g$'s. From $x_0=0.5$, $g_1$ diverges (then becomes undefined), and $g_2$ converges to $0.57$. D8.4 explains why.}
\end{deriv}

\begin{deriv}{D8.4}{Why $|g'|<1$ makes fixed point iteration converge}{\sone\ explains the fixed point theorem and linear convergence}
\goal{show that if $|g'(x)|\le k<1$, the errors $|x_n-p|$ shrink to 0.}
\idea{compare each iterate's distance to the fixed point $p$ with the previous one, using the Mean Value Theorem.}
\step{The fixed point satisfies $p=g(p)$. Subtract this from $x_{n+1}=g(x_n)$:}
\[x_{n+1}-p=g(x_n)-g(p)\]
\step{Mean Value Theorem: $g(x_n)-g(p)=g'(\xi_n)\,(x_n-p)$ for some $\xi_n$ between $x_n$ and $p$:}
\[x_{n+1}-p=g'(\xi_n)\,(x_n-p)\]
\step{Take absolute values, and use $|g'|\le k$:}
\[|x_{n+1}-p|\le k\,|x_n-p|\]
\step{Apply this repeatedly:}
\[|x_n-p|\le k\,|x_{n-1}-p|\le k^2|x_{n-2}-p|\le\cdots\le k^n|x_0-p|\ \longrightarrow\ 0\quad\text{since }0\le k<1\]
\result{the error shrinks by a factor of about $\lambda=|g'(p)|$ each step: \textbf{linear} convergence. If $|g'(p)|>1$ the error \emph{grows}.}
\end{deriv}

\begin{deriv}{D9.1}{Newton's method from the tangent line}{\sthree\ 2025 Q10 (3 marks)}
\goal{derive $x_{n+1}=x_n-\dfrac{f(x_n)}{f'(x_n)}$.}
\idea{replace the curve by its tangent line at $x_n$. Where the tangent hits zero is the next guess.}
\step{\mk{1} The tangent line to $f$ at $x=x_n$ (the degree-1 Taylor polynomial):}
\[f(x)\approx f(x_n)+f'(x_n)\,(x-x_n)\]
\step{\mk{1} Let $x_{n+1}$ be where the tangent line crosses the $x$-axis, i.e.\ set it equal to 0:}
\[0=f(x_n)+f'(x_n)\,(x_{n+1}-x_n)\]
\step{\mk{1} Solve for $x_{n+1}$: subtract $f(x_n)$, divide by $f'(x_n)$, then add $x_n$:}
\[x_{n+1}-x_n=-\frac{f(x_n)}{f'(x_n)}\quad\Rightarrow\quad x_{n+1}=x_n-\frac{f(x_n)}{f'(x_n)}\]
\result{Newton's method: $x_{n+1}=x_n-\dfrac{f(x_n)}{f'(x_n)}$. It needs $f'(x_n)\neq0$.}
\end{deriv}

\begin{deriv}{D9.2}{The secant method from Newton's method}{\stwo}
\goal{derive $x_{n+1}=x_n-f(x_n)\dfrac{x_n-x_{n-1}}{f(x_n)-f(x_{n-1})}$.}
\idea{Newton needs $f'(x_n)$. Replace it by the slope of the line through the last two points.}
\step{Approximate the derivative by the secant slope (rise over run between the last two iterates):}
\[f'(x_n)\approx\frac{f(x_n)-f(x_{n-1})}{x_n-x_{n-1}}\]
\step{Substitute into Newton. Dividing by a fraction means multiplying by its reciprocal:}
\[x_{n+1}=x_n-\frac{f(x_n)}{\dfrac{f(x_n)-f(x_{n-1})}{x_n-x_{n-1}}}=x_n-f(x_n)\,\frac{x_n-x_{n-1}}{f(x_n)-f(x_{n-1})}\]
\result{secant method. It needs two starting values and no derivative, and its order is about 1.618 (less than Newton's 2).}
\end{deriv}

\begin{deriv}{D9.3}{Iterations for $\sqrt a$, $\sqrt[3]a$ and more}{\sthree\ Q11: 2024 (Newton, $\sqrt[3]a$) and 2025 (secant, $\sqrt a$), 4 marks}
\goal{derive an iteration for a root of a number $a>0$.}
\idea{choose $f(x)$ whose root \emph{is} the number you want (for $\sqrt[k]a$, use $f(x)=x^k-a$), apply the method, then simplify.}

\textbf{(a) Heron's method: Newton for $\sqrt a$} (lecture)
\step{Choose $f(x)=x^2-a$ (its root is $\sqrt a$), so $f'(x)=2x$.}
\step{Newton, then split the fraction:}
\[x_{n+1}=x_n-\frac{x_n^2-a}{2x_n}=x_n-\frac{x_n}{2}+\frac{a}{2x_n}=\frac12\Big(x_n+\frac a{x_n}\Big)\]

\textbf{(b) Newton for $\sqrt[3]a$} (2024 Q11)
\step{\mk{1.5} Choose $f(x)=x^3-a$ (its root is $\sqrt[3]a$), so $f'(x)=3x^2$.}
\step{\mk{1.5} Apply Newton:}
\[x_{n+1}=x_n-\frac{x_n^3-a}{3x_n^2}\]
\step{\mk{1} Split the fraction into two pieces and simplify:}
\[x_{n+1}=x_n-\frac{x_n^3}{3x_n^2}+\frac{a}{3x_n^2}=x_n-\frac{x_n}{3}+\frac{a}{3x_n^2}=\frac{2x_n}{3}+\frac{a}{3x_n^2}\]

\textbf{(c) Secant for $\sqrt a$} (2025 Q11)
\step{\mk{1} Choose $f(x)=x^2-a$.}
\step{\mk{2} Apply the secant method. In the denominator the $a$'s cancel, leaving a difference of squares:}
\[x_{n+1}=x_n-(x_n^2-a)\,\frac{x_n-x_{n-1}}{(x_n^2-a)-(x_{n-1}^2-a)}=x_n-(x_n^2-a)\,\frac{x_n-x_{n-1}}{x_n^2-x_{n-1}^2}\]
\step{\mk{1} Factor $x_n^2-x_{n-1}^2=(x_n-x_{n-1})(x_n+x_{n-1})$ and cancel $(x_n-x_{n-1})$:}
\[x_{n+1}=x_n-\frac{x_n^2-a}{x_n+x_{n-1}}=x_n+\frac{a-x_n^2}{x_n+x_{n-1}}\]

\textbf{(d) Secant for $\sqrt[3]a$} (Practice Exam 2)
\step{Same as (c) with $f(x)=x^3-a$, using the difference of cubes $x_n^3-x_{n-1}^3=(x_n-x_{n-1})(x_n^2+x_nx_{n-1}+x_{n-1}^2)$:}
\[x_{n+1}=x_n-\frac{x_n^3-a}{x_n^2+x_nx_{n-1}+x_{n-1}^2}\]

\textbf{(e) Newton for $\sqrt[4]a$} (Practice Exam 1)
\step{$f(x)=x^4-a$, $f'(x)=4x^3$:\quad $x_{n+1}=x_n-\dfrac{x_n^4-a}{4x_n^3}=\dfrac{3x_n}{4}+\dfrac{a}{4x_n^3}$.}

\result{the recipe is always: $f=x^k-a$ $\to$ apply the method $\to$ split the fraction (Newton) or factor and cancel (secant).}
\end{deriv}

\begin{deriv}{D9.4}{Why Newton converges quadratically}{\sone\ explains ``the correct digits double''}
\goal{show $x_{n+1}-p\approx C\,(x_n-p)^2$, where $p$ is the root.}
\idea{Taylor-expand $f$ about $x_n$ and evaluate it at the root, where $f(p)=0$.}
\step{Taylor degree 1 with remainder, centred on $x_n$, evaluated at the root $p$ (where $f(p)=0$):}
\[0=f(p)=f(x_n)+f'(x_n)(p-x_n)+\frac{f''(\xi)}{2}(p-x_n)^2\]
\step{Divide by $f'(x_n)$ and rearrange:}
\[\underbrace{x_n-\frac{f(x_n)}{f'(x_n)}}_{=\,x_{n+1}}-\,p=\frac{f''(\xi)}{2f'(x_n)}\,(x_n-p)^2\]
\step{So the new error is (a constant) $\times$ (the old error)$^2$: if the error is $10^{-3}$, the next is roughly $10^{-6}$.}
\result{$|x_{n+1}-p|\approx\Big|\dfrac{f''(p)}{2f'(p)}\Big|\,|x_n-p|^2$: quadratic convergence. It breaks down if $f'(p)=0$ (a repeated root), where Newton is only linear.}
\end{deriv}

% ======================= WEEKS 10-11 =======================
\topic{Weeks 10--11: Numerical differentiation and integration}

\textbf{Notation.} $f_0=f(x_0)$, $f_0'=f'(x_0)$, $f_0''=f''(x_0)$, and so on. All the differentiation formulas come from Toolkit T1 (Taylor with step $h$).

\begin{deriv}{D10.1}{Forward and backward differences, with error terms}{\stwo}
\goal{derive $f'(x_0)=\dfrac{f(x_0+h)-f(x_0)}{h}-\dfrac h2f''(c)$, and the backward version.}
\idea{a degree-1 Taylor expansion contains $f'$. Solve for it.}
\step{Taylor degree 1 with remainder (T1), $c$ between $x_0$ and $x_0+h$:}
\[f(x_0+h)=f_0+h\,f_0'+\frac{h^2}{2}f''(c)\]
\step{Move $f_0$ to the left:}
\[f(x_0+h)-f_0=h\,f_0'+\frac{h^2}{2}f''(c)\]
\step{Divide by $h$:}
\[\frac{f(x_0+h)-f_0}{h}=f_0'+\frac h2f''(c)\]
\step{Rearrange for $f_0'$:}
\[f'(x_0)=\frac{f(x_0+h)-f(x_0)}{h}-\frac h2f''(c)\qquad\text{(forward, error }O(h))\]
\step{\textbf{Backward:} replace $h$ by $-h$ everywhere, then tidy the signs (multiply the top and bottom of the fraction by $-1$):}
\[f'(x_0)=\frac{f(x_0-h)-f(x_0)}{-h}-\frac{(-h)}{2}f''(c)=\frac{f(x_0)-f(x_0-h)}{h}+\frac h2f''(c)\]
\result{forward: $\dfrac{f(x_0+h)-f(x_0)}{h}$; \ backward: $\dfrac{f(x_0)-f(x_0-h)}{h}$. Both are first order, $O(h)$, with errors of opposite sign.}
\end{deriv}

\begin{deriv}{D10.2}{Central difference for $f'$, with error term}{\stwo}
\goal{derive $f'(x_0)=\dfrac{f(x_0+h)-f(x_0-h)}{2h}-\dfrac{h^2}{6}f'''(c)$.}
\idea{expand $f(x_0+h)$ and $f(x_0-h)$, then \textbf{subtract}: the even-power terms ($f_0$, $f_0''$) cancel, leaving $f_0'$.}
\step{Degree-2 expansions with remainders (T1), $c_1\in(x_0,x_0+h)$, $c_2\in(x_0-h,x_0)$:}
\begin{align*}
f(x_0+h)&=f_0+hf_0'+\frac{h^2}{2}f_0''+\frac{h^3}{6}f'''(c_1)\qquad(1)\\
f(x_0-h)&=f_0-hf_0'+\frac{h^2}{2}f_0''-\frac{h^3}{6}f'''(c_2)\qquad(2)
\end{align*}
\step{Compute (1) $-$ (2). The $f_0$ and $\frac{h^2}{2}f_0''$ terms cancel:}
\[f(x_0+h)-f(x_0-h)=2hf_0'+\frac{h^3}{6}\big[f'''(c_1)+f'''(c_2)\big]\]
\step{Divide by $2h$ and rearrange for $f_0'$:}
\[f'(x_0)=\frac{f(x_0+h)-f(x_0-h)}{2h}-\frac{h^2}{6}\cdot\frac{f'''(c_1)+f'''(c_2)}{2}\]
\step{Merge the two unknown points with the IVT trick (T3): the average equals $f'''(c)$ for some $c$ between $c_1$ and $c_2$.}
\result{$f'(x_0)=\dfrac{f(x_0+h)-f(x_0-h)}{2h}-\dfrac{h^2}{6}f'''(c)$, \ $c\in(x_0-h,x_0+h)$. Second order, $O(h^2)$.}
\end{deriv}

\begin{deriv}{D10.3}{Central difference for $f''$, with error term}{\sthree\ 2023 Q11a (4 marks)}
\goal{derive $f''(x_0)=\dfrac{f(x_0-h)-2f(x_0)+f(x_0+h)}{h^2}-\dfrac{h^2}{12}f^{(4)}(c)$, \ $c\in(x_0-h,x_0+h)$.}
\idea{expand $f(x_0\pm h)$ up to $h^3$, then \textbf{add}: the odd-power terms ($f_0'$, $f_0'''$) cancel, leaving $f_0''$.}
\step{\mk{1} Degree-3 expansions with $h^4$ remainders, $c_1\in(x_0,x_0+h)$, $c_2\in(x_0-h,x_0)$:}
\begin{align*}
f(x_0+h)&=f_0+hf_0'+\frac{h^2}{2}f_0''+\frac{h^3}{6}f_0'''+\frac{h^4}{24}f^{(4)}(c_1)\qquad(1)\\
f(x_0-h)&=f_0-hf_0'+\frac{h^2}{2}f_0''-\frac{h^3}{6}f_0'''+\frac{h^4}{24}f^{(4)}(c_2)\qquad(2)
\end{align*}
\why{we go up to $h^3$ because, after adding and dividing by $h^2$, we need to see the size of the error.}
\step{\mk{1} Add (1) $+$ (2). The $\pm hf_0'$ and $\pm\frac{h^3}{6}f_0'''$ terms cancel:}
\[f(x_0+h)+f(x_0-h)=2f_0+h^2f_0''+\frac{h^4}{24}\big[f^{(4)}(c_1)+f^{(4)}(c_2)\big]\]
\step{\mk{1} Rearrange for $f_0''$: subtract $2f_0$, then divide by $h^2$:}
\[f''(x_0)=\frac{f(x_0-h)-2f(x_0)+f(x_0+h)}{h^2}-\frac{h^2}{24}\big[f^{(4)}(c_1)+f^{(4)}(c_2)\big]\]
\[\phantom{f''(x_0)}=\frac{f(x_0-h)-2f(x_0)+f(x_0+h)}{h^2}-\frac{h^2}{12}\cdot\frac{f^{(4)}(c_1)+f^{(4)}(c_2)}{2}\]
\step{\mk{1} IVT (T3): provided $f^{(4)}$ is continuous on $[x_0-h,x_0+h]$, there is a $c$ between $c_1$ and $c_2$ with $f^{(4)}(c)=\dfrac{f^{(4)}(c_1)+f^{(4)}(c_2)}{2}$.}
\result{$f''(x_0)=\dfrac{f(x_0-h)-2f(x_0)+f(x_0+h)}{h^2}-\dfrac{h^2}{12}f^{(4)}(c)$. Second order, $O(h^2)$.}
\trap{forgetting to \emph{state} the IVT step. It is worth a mark on its own.}
\end{deriv}

\begin{deriv}{D10.4}{Second-order forward difference}{\stwo\ lecture problem}
\goal{derive $f'(x_0)\approx\dfrac{-3f(x_0)+4f(x_0+h)-f(x_0+2h)}{2h}$ with error $O(h^2)$.}
\idea{use two points to the right. Combine their expansions so the $f''$ terms cancel.}
\step{Expand both up to $h^3$ (for $x_0+2h$, put $2h$ in place of $h$: $\frac{(2h)^2}{2}=2h^2$ and $\frac{(2h)^3}{6}=\frac43h^3$):}
\begin{align*}
f(x_0+h)&=f_0+hf_0'+\frac{h^2}{2}f_0''+\frac{h^3}{6}f_0'''+O(h^4)\qquad(1)\\
f(x_0+2h)&=f_0+2hf_0'+2h^2f_0''+\frac{4h^3}{3}f_0'''+O(h^4)\qquad(2)
\end{align*}
\step{Choose a combination that kills $f_0''$: (1) has $\frac{h^2}{2}f_0''$ and (2) has $2h^2f_0''$, so take $4\times(1)-(2)$:}
\[4f(x_0+h)-f(x_0+2h)=3f_0+2hf_0'+\Big(\frac{4}{6}-\frac43\Big)h^3f_0'''+O(h^4)=3f_0+2hf_0'-\frac23h^3f_0'''+O(h^4)\]
\step{Rearrange for $f_0'$: move $3f_0$ across, then divide by $2h$:}
\[f'(x_0)=\frac{-3f(x_0)+4f(x_0+h)-f(x_0+2h)}{2h}+\frac{h^2}{3}f_0'''+O(h^3)\]
\result{second-order accuracy using only points to the right (useful at the left end of a table).}
\end{deriv}

\begin{deriv}{D10.5}{How the choice of $h$ affects accuracy}{\sthree\ 2023 Q11b (2 marks) \quad+\quad \sone\ the calculation}
\goal{explain why $h$ can be too large \emph{or} too small.}
\step{\mk{1 + 1} \textbf{Exam answer:} ``Too large a value of $h$ gives inaccuracy from \textbf{truncation error} (the Taylor terms we dropped, $O(h^p)$). Too small a value gives inaccuracy from \textbf{roundoff error}: the numerator subtracts two nearly-equal numbers (catastrophic cancellation), and dividing by a tiny $h$ magnifies it.''}

\step{\textbf{(Understanding) Model the forward difference's total error.} Each stored value of $f$ has relative rounding error up to $\varepsilon$, so the numerator can be off by about $2\varepsilon|f|$. After dividing by $h$:}
\[E(h)\approx\underbrace{\frac M2\,h}_{\text{truncation}}+\underbrace{\frac{2\varepsilon|f|}{h}}_{\text{roundoff}}\qquad(M\approx|f''|)\]
\step{Minimise: set $E'(h)=\frac M2-\frac{2\varepsilon|f|}{h^2}=0$, giving $h=2\sqrt{\varepsilon|f|/M}$. With $|f|\approx M$, this is about $\sqrt\varepsilon\approx10^{-8}$ in double precision.}
\result{for the forward difference the best $h$ is about $\sqrt{\varepsilon}\approx10^{-8}$ (the bottom of the V-shaped error plot). For the central difference it is about $\varepsilon^{1/3}\approx10^{-5}$.}
\end{deriv}

\begin{deriv}{D11.1}{Trapezoidal rule for one strip}{\sthree\ 2023 Q13a (3 marks)}
\goal{show $\displaystyle\int_{x_0}^{x_1}P_1(x)\,dx=\frac h2\big[f(x_0)+f(x_1)\big]$, where $P_1$ interpolates $f$ at $x_0$ and $x_1$.}
\idea{write the straight line in Newton forward difference form, switch to $s$ (T4), then integrate a simple polynomial in $s$.}
\step{\mk{1} The degree-1 interpolant in forward difference form:}
\[P_1(x)=y_0+s\,\Delta y_0,\qquad s=\frac{x-x_0}{h},\qquad\Delta y_0=y_1-y_0\]
\step{Change variable (T4): $x=x_0+sh$ gives $dx=h\,ds$. At $x=x_0$, $s=0$; at $x=x_1$, $s=1$.}
\step{\mk{1} Integrate in $s$:}
\[\int_{x_0}^{x_1}P_1\,dx=\int_0^1\big(y_0+s\,\Delta y_0\big)\,h\,ds=h\Big[y_0s+\tfrac12s^2\Delta y_0\Big]_0^1=h\Big(y_0+\tfrac12\Delta y_0\Big)\]
\step{\mk{1} Substitute $\Delta y_0=y_1-y_0$ and simplify:}
\[h\Big(y_0+\tfrac12(y_1-y_0)\Big)=h\Big(\tfrac12y_0+\tfrac12y_1\Big)=\frac h2\big[f(x_0)+f(x_1)\big]\]
\result{$\displaystyle\int_{x_0}^{x_1}f\,dx\approx\frac h2\big[f(x_0)+f(x_1)\big]$, which is the area of a trapezoid.}
\end{deriv}

\begin{deriv}{D11.2}{Composite trapezoidal rule}{\sthree\ 2023 Q13b (2 marks)}
\goal{show $\displaystyle\int_a^bf\,dx\approx\frac h2\Big[f(x_0)+2\sum_{i=1}^{n-1}f(x_i)+f(x_n)\Big]$, with $h=\frac{b-a}{n}$.}
\idea{split $[a,b]$ into $n$ strips, apply D11.1 to each, and add.}
\step{\mk{1} Split the integral into strips:}
\[\int_a^bf\,dx=\int_{x_0}^{x_1}f\,dx+\int_{x_1}^{x_2}f\,dx+\cdots+\int_{x_{n-1}}^{x_n}f\,dx\]
\step{\mk{0.5} Apply D11.1 to every strip:}
\[\approx\frac h2\big[f_0+f_1\big]+\frac h2\big[f_1+f_2\big]+\cdots+\frac h2\big[f_{n-1}+f_n\big]\]
\step{\mk{0.5} Collect terms. Each interior point $x_1,\dots,x_{n-1}$ is the right end of one strip and the left end of the next, so it appears \textbf{twice}. The endpoints appear once:}
\[=\frac h2\Big[f(x_0)+2\big(f(x_1)+\cdots+f(x_{n-1})\big)+f(x_n)\Big]\]
\result{composite trapezoidal rule, with weights $1,2,2,\dots,2,1$. Error $-\frac{(b-a)h^2}{12}f''(c)$, so second order.}
\end{deriv}

\begin{deriv}{D11.3}{Simpson's rule for one pair of strips}{\sthree\ 2024 Q13 (5 marks)}
\goal{show $\displaystyle\int_{x_0}^{x_2}P_2(x)\,dx=\frac h3\big[f(x_0)+4f(x_1)+f(x_2)\big]$.}
\idea{same as D11.1, but with the degree-2 interpolant and $s$ running from 0 to 2.}
\step{\mk{1} The degree-2 interpolant in forward difference form. Expand $\frac{s(s-1)}2=\frac12s^2-\frac12s$ ready for integrating:}
\[P_2(x)=y_0+s\,\Delta y_0+\Big(\tfrac12s^2-\tfrac12s\Big)\Delta^2y_0\]
\step{Change variable: $dx=h\,ds$. At $x=x_0$, $s=0$; at $x=x_2$, $s=2$.}
\step{\mk{1} Integrate term by term:}
\[\int_{x_0}^{x_2}P_2\,dx=h\Big[y_0s+\tfrac12s^2\Delta y_0+\Big(\tfrac16s^3-\tfrac14s^2\Big)\Delta^2y_0\Big]_0^2\]
\step{Evaluate at $s=2$ (everything is 0 at $s=0$). Note that $\frac16(8)-\frac14(4)=\frac43-1=\frac13$:}
\[=h\Big[2y_0+2\Delta y_0+\tfrac13\Delta^2y_0\Big]\]
\step{\mk{1} Substitute $\Delta y_0=y_1-y_0$ and $\Delta^2y_0=y_2-2y_1+y_0$:}
\[=h\Big[2y_0+2y_1-2y_0+\tfrac13y_2-\tfrac23y_1+\tfrac13y_0\Big]\]
\step{\mk{1} Collect each $y$: $y_0$: $2-2+\frac13=\frac13$; \ $y_1$: $2-\frac23=\frac43$; \ $y_2$: $\frac13$:}
\[=h\Big[\tfrac13y_0+\tfrac43y_1+\tfrac13y_2\Big]\]
\step{\mk{1} Factor out $\frac13$:}
\[=\frac h3\big[f(x_0)+4f(x_1)+f(x_2)\big]\]
\result{$\displaystyle\int_{x_0}^{x_2}f\,dx\approx\frac h3\big[f(x_0)+4f(x_1)+f(x_2)\big]$, the area under the parabola through the three points.}
\end{deriv}

\begin{deriv}{D11.4}{Composite Simpson's rule}{\stwo}
\goal{show $\displaystyle\int_a^bf\,dx\approx\frac h3\Big[f_0+4\!\!\sum_{\text{odd }i}\!\!f_i+2\!\!\!\sum_{\text{even interior }i}\!\!\!f_i+f_n\Big]$, with $n$ even.}
\idea{apply D11.3 to consecutive \emph{pairs} of strips, then add. That's why $n$ must be even.}
\step{Split into pairs: $[x_0,x_2],[x_2,x_4],\dots,[x_{n-2},x_n]$. There are $\frac n2$ pairs, so $n$ must be even.}
\step{Apply D11.3 to each pair:}
\[\frac h3\big[f_0+4f_1+f_2\big]+\frac h3\big[f_2+4f_3+f_4\big]+\cdots+\frac h3\big[f_{n-2}+4f_{n-1}+f_n\big]\]
\step{Collect terms. The \textbf{odd} points are pair middles, so they get weight 4. Each \textbf{even interior} point is shared by two pairs ($1+1$), so it gets weight 2. The ends get weight 1.}
\result{weights $1,4,2,4,\dots,2,4,1$, multiplied by $\frac h3$. Error $-\frac{(b-a)h^4}{180}f^{(4)}(c)$, so fourth order.}
\end{deriv}

\begin{deriv}{D11.5}{Why Trapezoid is exact for lines and Simpson is exact for cubics}{\sone}
\goal{see where the ``exact for degree $\le1$'' and ``exact for degree $\le3$'' facts come from.}
\step{\textbf{From the error terms.} The trapezoid's error contains $f''$, and $f''=0$ for any straight line, so the error is 0. Simpson's error contains $f^{(4)}$, and $f^{(4)}=0$ for any cubic.}
\step{\textbf{A direct check that Simpson is exact for $x^3$.} On $[-h,h]$ (nodes $-h,0,h$):}
\[\int_{-h}^{h}x^3\,dx=0\ \ (\text{odd function}),\qquad\frac h3\big[(-h)^3+4\cdot0+h^3\big]=0\ \cmark\]
\step{\textbf{And for $x^2$:}}
\[\int_{-h}^{h}x^2\,dx=\frac{2h^3}{3},\qquad\frac h3\big[h^2+4\cdot0+h^2\big]=\frac{2h^3}{3}\ \cmark\]
\result{Simpson fits a parabola but is exact one degree higher, for cubics: the symmetric weights cancel the $x^3$ error.}
\end{deriv}

% ======================= WEEKS 12-13 =======================
\topic{Weeks 12--13: Linear systems}

\begin{deriv}{D12.1}{Forward substitution}{\stwo}
\goal{solve a lower triangular system $A\mathbf x=\mathbf b$ (all zeros above the diagonal).}
\idea{row 1 contains only $x_1$. Once $x_1$ is known, row 2 contains only one new unknown, and so on downwards.}
\step{Row $i$ of a lower triangular system only involves $x_1,\dots,x_i$ (the entries $a_{ij}$ with $j>i$ are 0):}
\[a_{i1}x_1+a_{i2}x_2+\cdots+a_{i,i-1}x_{i-1}+a_{ii}x_i=b_i\]
\step{Move the already-known terms to the right-hand side, then divide by the diagonal entry:}
\[x_i=\frac{b_i-\sum_{j=1}^{i-1}a_{ij}x_j}{a_{ii}},\qquad i=1,2,\dots,n\ \ \text{(top to bottom)}\]
\result{forward substitution. It needs every $a_{ii}\neq0$.}
\end{deriv}

\begin{deriv}{D12.2}{Backward substitution}{\stwo}
\goal{solve an upper triangular system (all zeros below the diagonal).}
\idea{the last row contains only $x_n$. Work upwards.}
\step{Row $i$ only involves $x_i,\dots,x_n$:\quad $a_{ii}x_i+a_{i,i+1}x_{i+1}+\cdots+a_{in}x_n=b_i$.}
\step{Move the known terms (those to the right of the diagonal) across, then divide:}
\[x_i=\frac{b_i-\sum_{j=i+1}^{n}a_{ij}x_j}{a_{ii}},\qquad i=n,n-1,\dots,1\ \ \text{(bottom to top)}\]
\result{backward substitution. It needs every $a_{ii}\neq0$.}
\end{deriv}

\begin{deriv}{D12.3}{How LU factorisation solves $A\mathbf x=\mathbf b$}{\sthree\ 2023 Q15b (3 marks)}
\goal{explain, in words and symbols, how $A=LU$ is used to solve $A\mathbf x=\mathbf b$.}
\idea{one hard system becomes two easy triangular ones.}
\step{\mk{1} Replace $A$ by its factorisation, and regroup the brackets (matrix multiplication is associative):}
\[A\mathbf x=\mathbf b\ \Rightarrow\ (LU)\mathbf x=\mathbf b\ \Rightarrow\ L\,(U\mathbf x)=\mathbf b\]
\step{Name the bracket $\mathbf z=U\mathbf x$. Then the system reads $L\mathbf z=\mathbf b$.}
\step{\mk{1} Solve $L\mathbf z=\mathbf b$ for $\mathbf z$ by \textbf{forward substitution} ($L$ is lower triangular).}
\step{\mk{1} Solve $U\mathbf x=\mathbf z$ for $\mathbf x$ by \textbf{backward substitution} ($U$ is upper triangular).}
\result{$A\mathbf x=\mathbf b$ $\to$ $L\mathbf z=\mathbf b$ (forward) $\to$ $U\mathbf x=\mathbf z$ (backward). For a new $\mathbf b$, reuse the same $L$ and $U$.}
\end{deriv}

\begin{deriv}{D12.4}{Doolittle LU formulas (3\,$\times$\,3)}{\sthree\ Q15 every year (2--3 marks for the missing entries)}
\goal{find formulas for the entries of $L$ (with 1s on its diagonal) and $U$ such that $A=LU$.}
\idea{multiply $L$ and $U$ out symbolically, set each entry equal to the matching entry of $A$, and solve in an order where only one entry is unknown at a time.}
\step{Write out the factors:}
\[L=\begin{bmatrix}1&0&0\\\ell_{21}&1&0\\\ell_{31}&\ell_{32}&1\end{bmatrix},\qquad U=\begin{bmatrix}u_{11}&u_{12}&u_{13}\\0&u_{22}&u_{23}\\0&0&u_{33}\end{bmatrix}\]
\step{Multiply them: entry $(i,j)$ is (row $i$ of $L$) $\cdot$ (column $j$ of $U$):}
\[LU=\begin{bmatrix}u_{11}&u_{12}&u_{13}\\[2pt]\ell_{21}u_{11}&\ell_{21}u_{12}+u_{22}&\ell_{21}u_{13}+u_{23}\\[2pt]\ell_{31}u_{11}&\ell_{31}u_{12}+\ell_{32}u_{22}&\ell_{31}u_{13}+\ell_{32}u_{23}+u_{33}\end{bmatrix}\]
\step{Set $LU=A$ entry by entry and solve in this order: \textbf{row 1 of $U$}, then \textbf{column 1 of $L$}, then \textbf{row 2 of $U$}, then \textbf{column 2 of $L$}, then \textbf{$u_{33}$}:}
\begin{align*}
&\text{row 1 of }U: && u_{11}=a_{11},\quad u_{12}=a_{12},\quad u_{13}=a_{13}\\
&\text{column 1 of }L: && \ell_{21}=\frac{a_{21}}{u_{11}},\quad \ell_{31}=\frac{a_{31}}{u_{11}}\\
&\text{row 2 of }U: && u_{22}=a_{22}-\ell_{21}u_{12},\quad u_{23}=a_{23}-\ell_{21}u_{13}\\
&\text{column 2 of }L: && \ell_{32}=\frac{a_{32}-\ell_{31}u_{12}}{u_{22}}\\
&\text{last entry}: && u_{33}=a_{33}-\ell_{31}u_{13}-\ell_{32}u_{23}
\end{align*}
\why{in this order, each equation contains exactly one entry you haven't found yet.}

\textbf{Worked example (2024 Q15a).} \emph{What the question gives you:} $A$, and a partly filled-in $L$ and $U$. The boxed entries are the blanks:
\[\underbrace{\begin{bmatrix}2&1&-4\\-4&-5&10\\6&6&-10\end{bmatrix}}_{A}=\underbrace{\begin{bmatrix}1&0&0\\-2&1&0\\3&\boxed{\ell_{32}}&1\end{bmatrix}}_{L}\underbrace{\begin{bmatrix}2&1&-4\\0&\boxed{u_{22}}&2\\0&0&\boxed{u_{33}}\end{bmatrix}}_{U}\]
\emph{Your job:} find the three boxed numbers. For each one, take the row of $L$ and column of $U$ that contain it, set their dot product equal to the matching entry of $A$, and solve. (The formulas from Step 3 give the same thing.)

\step{\textbf{Find $u_{22}$} using row 2 of $L$ and column 2 of $U$, which must give $a_{22}=-5$:}
\[(-2,\ 1,\ 0)\cdot(1,\ u_{22},\ 0)=-2+u_{22}=-5\ \Rightarrow\ u_{22}=-3\qquad\Big(\text{formula: }a_{22}-\ell_{21}u_{12}=-5-(-2)(1)\Big)\]

\step{\textbf{Find $\ell_{32}$} using row 3 of $L$ and column 2 of $U$ (with $u_{22}=-3$ now filled in), which must give $a_{32}=6$:}
\[(3,\ \ell_{32},\ 1)\cdot(1,\ -3,\ 0)=3-3\ell_{32}=6\ \Rightarrow\ \ell_{32}=-1\qquad\Big(\text{formula: }\tfrac{a_{32}-\ell_{31}u_{12}}{u_{22}}=\tfrac{6-3}{-3}\Big)\]

\step{\textbf{Find $u_{33}$} using row 3 of $L$ (with $\ell_{32}=-1$) and column 3 of $U$, which must give $a_{33}=-10$:}
\[(3,\ -1,\ 1)\cdot(-4,\ 2,\ u_{33})=-12-2+u_{33}=-10\ \Rightarrow\ u_{33}=4\qquad\Big(\text{formula: }a_{33}-\ell_{31}u_{13}-\ell_{32}u_{23}\Big)\]

\chk{row 3 of $L$ times $U$: $\big(3\cdot2,\ 3\cdot1+(-1)(-3),\ -10\big)=(6,6,-10)$, which is row 3 of $A$ \cmark}
\why{order matters: $u_{22}$ must come before $\ell_{32}$, because the $\ell_{32}$ equation uses $u_{22}$.}
\result{every Doolittle formula comes from one rule: $a_{ij}=(\text{row }i\text{ of }L)\cdot(\text{column }j\text{ of }U)$.}
\trap{a zero pivot ($u_{11}=0$ or $u_{22}=0$) makes a division impossible. Pivoting (swapping rows) fixes it.}
\end{deriv}

\begin{deriv}{D13.1}{Showing a matrix is positive definite}{\sthree\ 2025 Q16a (3 marks)}
\goal{show $A=\begin{bmatrix}1&-1\\-1&2\end{bmatrix}$ is positive definite, i.e.\ $\mathbf x^TA\mathbf x>0$ for every $\mathbf x\neq\mathbf0$.}
\idea{multiply $\mathbf x^TA\mathbf x$ out into a quadratic, then complete the square to show it's a sum of squares.}
\step{\mk{1} State what has to be shown: $\mathbf x^TA\mathbf x>0$ for all nonzero $\mathbf x=(x_1,x_2)^T$.}
\step{\mk{1} Multiply out, starting with $A\mathbf x$:}
\[A\mathbf x=\begin{bmatrix}x_1-x_2\\-x_1+2x_2\end{bmatrix},\qquad\mathbf x^TA\mathbf x=x_1(x_1-x_2)+x_2(-x_1+2x_2)=x_1^2-2x_1x_2+2x_2^2\]
\step{\mk{0.5} Complete the square. Split $2x_2^2=x_2^2+x_2^2$ so that a perfect square appears:}
\[x_1^2-2x_1x_2+x_2^2+x_2^2=(x_1-x_2)^2+x_2^2\]
\step{\mk{0.5} Conclude. A sum of squares is $\ge0$. It equals 0 only if $x_2=0$ \emph{and} $x_1=x_2$, which means $\mathbf x=\mathbf0$. So $\mathbf x^TA\mathbf x>0$ for every nonzero $\mathbf x$, and $A$ is positive definite.}

\textbf{General 2\,$\times$\,2 version} \sone. For $A=\begin{bmatrix}a&b\\b&c\end{bmatrix}$ with $a>0$:
\[\mathbf x^TA\mathbf x=ax_1^2+2bx_1x_2+cx_2^2=a\Big(x_1+\frac bax_2\Big)^2+\Big(c-\frac{b^2}{a}\Big)x_2^2\]
so $A$ is positive definite exactly when $a>0$ and $ac-b^2>0$.
\chk{for $\begin{bmatrix}1&-1\\-1&2\end{bmatrix}$: $a=1>0$ and $ac-b^2=2-1=1>0$ \cmark}
\result{to prove positive definiteness: multiply out $\mathbf x^TA\mathbf x$, write it as a sum of squares, and argue it is 0 only for $\mathbf x=\mathbf0$.}
\end{deriv}

\begin{deriv}{D13.2}{Cholesky formulas (3\,$\times$\,3)}{\sthree\ Q16 every year (2--3 marks for the missing entries)}
\goal{find the entries of the lower triangular $L$ with $A=LL^T$, for a symmetric positive definite $A$.}
\idea{exactly like D12.4: multiply out, equate entries with $A$, and solve column by column.}
\step{Write the factors ($L^T$ is $L$ reflected across its diagonal):}
\[L=\begin{bmatrix}\ell_{11}&0&0\\\ell_{21}&\ell_{22}&0\\\ell_{31}&\ell_{32}&\ell_{33}\end{bmatrix},\qquad L^T=\begin{bmatrix}\ell_{11}&\ell_{21}&\ell_{31}\\0&\ell_{22}&\ell_{32}\\0&0&\ell_{33}\end{bmatrix}\]
\step{Multiply. Only the lower triangle is needed, because $A$ is symmetric:}
\[LL^T=\begin{bmatrix}\ell_{11}^2&\cdot&\cdot\\[2pt]\ell_{21}\ell_{11}&\ell_{21}^2+\ell_{22}^2&\cdot\\[2pt]\ell_{31}\ell_{11}&\ell_{31}\ell_{21}+\ell_{32}\ell_{22}&\ell_{31}^2+\ell_{32}^2+\ell_{33}^2\end{bmatrix}\]
\step{Set this equal to $A$ and solve \textbf{column by column}:}
\begin{align*}
&\text{column 1}: && \ell_{11}=\sqrt{a_{11}},\quad \ell_{21}=\frac{a_{21}}{\ell_{11}},\quad \ell_{31}=\frac{a_{31}}{\ell_{11}}\\
&\text{column 2}: && \ell_{22}=\sqrt{a_{22}-\ell_{21}^2},\quad \ell_{32}=\frac{a_{32}-\ell_{31}\ell_{21}}{\ell_{22}}\\
&\text{column 3}: && \ell_{33}=\sqrt{a_{33}-\ell_{31}^2-\ell_{32}^2}
\end{align*}
\why{a diagonal entry comes from a sum of squares, so it needs a square root (take the positive one). An entry below the diagonal comes from a dot product, so it needs a division.}

\textbf{Worked example (2025 Q16b).} \emph{What the question gives you:} $A$, and a partly filled-in $L$ (with $L^T$ beside it). The boxed entries are the blanks:
\[\underbrace{\begin{bmatrix}36&-12&-6\\-12&13&11\\-6&11&35\end{bmatrix}}_{A}=\underbrace{\begin{bmatrix}\boxed{\ell_{11}}&0&0\\-2&3&0\\-1&\boxed{\ell_{32}}&\boxed{\ell_{33}}\end{bmatrix}}_{L}\underbrace{\begin{bmatrix}\ell_{11}&-2&-1\\0&3&\ell_{32}\\0&0&\ell_{33}\end{bmatrix}}_{L^T}\]
\emph{Your job:} find the three boxed numbers. Same idea as LU: (row of $L$) $\cdot$ (column of $L^T$) $=$ entry of $A$. Column $j$ of $L^T$ is just row $j$ of $L$, so entry $a_{ij}$ is (row $i$ of $L$) $\cdot$ (row $j$ of $L$).

\step{\textbf{Find $\ell_{11}$} using row 1 $\cdot$ row 1, which must give $a_{11}=36$:}
\[(\ell_{11},0,0)\cdot(\ell_{11},0,0)=\ell_{11}^2=36\ \Rightarrow\ \ell_{11}=6\qquad\big(\text{take the positive root}\big)\]

\step{\textbf{Find $\ell_{32}$} using row 3 $\cdot$ row 2, which must give $a_{32}=11$:}
\[(-1,\ \ell_{32},\ \ell_{33})\cdot(-2,\ 3,\ 0)=2+3\ell_{32}=11\ \Rightarrow\ \ell_{32}=3\qquad\Big(\text{formula: }\tfrac{a_{32}-\ell_{31}\ell_{21}}{\ell_{22}}=\tfrac{11-2}{3}\Big)\]

\step{\textbf{Find $\ell_{33}$} using row 3 $\cdot$ row 3, which must give $a_{33}=35$:}
\[(-1,\ 3,\ \ell_{33})\cdot(-1,\ 3,\ \ell_{33})=1+9+\ell_{33}^2=35\ \Rightarrow\ \ell_{33}^2=25\ \Rightarrow\ \ell_{33}=5\]

\chk{row 2 $\cdot$ row 2 $=(-2)^2+3^2=13=a_{22}$ \cmark, \quad row 2 $\cdot$ row 1 $=(-2)(6)=-12=a_{21}$ \cmark}
\result{diagonal: $\sqrt{\text{entry}-\text{sum of squares so far}}$; \ below the diagonal: $(\text{entry}-\text{dot product so far})\div\text{diagonal}$.}
\trap{a negative (or zero) number under a square root means $A$ is \textbf{not} positive definite. The factorisation fails, so switch to LU.}
\end{deriv}

\begin{deriv}{D13.3}{How Cholesky solves $A\mathbf x=\mathbf b$}{\stwo}
\goal{solve $A\mathbf x=\mathbf b$ using $A=LL^T$.}
\step{Replace $A$ and regroup: $\ (LL^T)\mathbf x=\mathbf b\ \Rightarrow\ L\,(L^T\mathbf x)=\mathbf b$.}
\step{Let $\mathbf z=L^T\mathbf x$. Solve $L\mathbf z=\mathbf b$ by \textbf{forward} substitution.}
\step{Solve $L^T\mathbf x=\mathbf z$ by \textbf{backward} substitution ($L^T$ is upper triangular).}
\result{the same two-step idea as LU, with $U=L^T$. It is about half the work, needs no pivoting, and succeeds exactly when $A$ is symmetric positive definite.}
\end{deriv}

\vspace{8pt}
\begin{tcolorbox}[enhanced,colback=accentlight,colframe=accent!50,boxrule=0.5pt,arc=2pt]
\textbf{Final tip.} Many of these derivations share a skeleton:
\begin{itemize}[leftmargin=16pt,itemsep=1pt]
\item ODE methods: Taylor about $t_i$ $\to$ evaluate at $t_{i+1}$ $\to$ $y'=f$ $\to$ drop $O(h^{\cdot})$.
\item Finite differences: expand $f(x_0\pm h)$ $\to$ add or subtract $\to$ rearrange $\to$ IVT.
\item Quadrature: interpolant in $s$ $\to$ $dx=h\,ds$ $\to$ integrate $\to$ substitute the $\Delta$'s.
\item Matrix factorisations: multiply out $\to$ equate entries $\to$ solve one unknown at a time.
\end{itemize}
Learn the skeletons, and the details become much easier to remember.
\end{tcolorbox}


\end{document}
